\chapter{Арифметика}\label{chap:arithmetic}\label{chap:arithmetics}

\section{Введение}

Цель этой главы — подготовить читателя, владеющего лишь базовой школьной алгеброй, к изучению арифметики. Мы начнём с краткого повторения основ целочисленной арифметики, обсудим деление столбиком (long division), наибольший общий делитель (greatest common divisor) и \concept{деление с остатком} (Euclidean division). После этого мы введём модульную арифметику (modular arithmetic) как \hilight{наиболее важный} навык для вычисления наших примеров «от руки». Затем мы познакомимся с многочленами (polynomials), вычислим их аналоги целочисленной арифметики и введём важное понятие \concept{интерполяции Лагранжа} (Lagrange interpolation).

\section{Целочисленная арифметика}
\label{integer_arithmetic}
В некотором смысле целочисленная арифметика лежит в основе значительной части современной криптографии. К счастью, большинство читателей, вероятно, помнят целочисленную арифметику из школы. Однако важно, чтобы вы могли уверенно применять эти понятия для понимания и выполнения вычислений во множестве примеров «от руки», составляющих неотъемлемую часть MoonMath Manual. Поэтому мы кратко повторим основные арифметические понятия, чтобы освежить вашу память и заполнить пробелы в знаниях.

Хотя термины и понятия этой главы могут непосредственно не встречаться в литературе по доказательствам с нулевым разглашением (zero-knowledge proofs), их понимание необходимо для изучения последующих глав и не только: такие термины, как \term{группы} (groups) или \term{поля} (fields), также очень часто встречаются в научных работах по криптографии с нулевым разглашением.

Многие идеи, представленные в этой главе, излагаются в рамках базового математического образования в большинстве школ по всему миру. Большая часть идей этого раздела может быть найдена в \cite{wu-1}. Более ориентированный на информатику подход представлен в \cite{mignotte-1992}.

\subsection{Целые числа, натуральные числа и рациональные числа}

Целые числа (integers) — это числа, которые можно записать без дробной части. Их также называют \term{целыми числами} (whole numbers).\footnote{Целые числа (whole numbers), а также их основные законы операций, вводятся, например, в главах 1–6 \cite{wu-1}.} Примерами чисел, которые \hilight{не} являются целыми, являются $\frac{2}{3}$, $1.2$ и $-1280.006$.

На протяжении всей книги мы используем символ $\mathbb{Z}$ как сокращение для множества всех \term{целых чисел}:

\begin{equation}
\label{integer_symbol}
\Z := \{\ldots, -3,-2,-1,0,1,2,3,\ldots\}
\end{equation}

Если $a\in \Z$ — целое число, то мы пишем $|a|$ для \term{абсолютной величины} (absolute value) $a$, то есть неотрицательного значения $a$ без учёта знака:

\begin{equation}
|4|= 4
\end{equation}
\begin{equation}
|-4|= 4
\end{equation}

Мы используем символ $\N$ для множества всех положительных целых чисел, обычно называемого множеством \term{натуральных чисел} (natural numbers). Кроме того, мы используем $\NN$ для множества всех неотрицательных целых чисел. Это означает, что $\mathbb{N}$ не содержит числа $0$, тогда как $\NN$ содержит:
\begin{align*}
\label{natural_symbol}
\N := \{1,2,3,\ldots\} & & \NN := \{0,1,2,3,\ldots\}
\end{align*}

Кроме того, мы используем символ $\Q$ для множества всех \term{рациональных чисел} (rational numbers), которые можно представить как множество всех дробей $\frac{n}{m}$, где $n\in \Z$ — целое число, а $m\in \N$ — натуральное число, таких, что не существует другой дроби $\frac{n'}{m'}$ и натурального числа $k\in \N$ с $k\neq 1$ и следующим уравнением:\footnote{Более глубокое введение в рациональные числа, их представление, а также арифметические операции над ними можно найти в части 2 \cite{wu-1} и в \chaptname{} 1, \secname{} 2 \cite{mignotte-1992}.}
\begin{equation}
\frac{n}{m} = \frac{k\cdot n'}{k\cdot m'}
\end{equation}

Множества $\N$, $\Z$ и $\Q$ имеют определённые на них понятия сложения и умножения. Большинство из нас, вероятно, способно выполнять многие вычисления с целыми числами в уме, но это становится всё труднее по мере увеличения их сложности. Мы будем часто использовать систему SageMath (\ref{sagemath_setup}) для более сложных вычислений (кольца и поля мы определим позже в этой книге):
\begin{sagecommandline}
sage: ZZ #  Sage notation for the set of integers
sage: NN # Sage notation for the set of natural numbers
sage: QQ # Sage notation for the set of rational numbers
sage: ZZ(5) # Get an element from the set of integers
sage: ZZ(5) + ZZ(3)
sage: ZZ(5) * NN(3)
sage: ZZ.random_element(10**6)
\end{sagecommandline}

Нас особенно интересует множество \term{простых чисел} (prime numbers). Простое число — это натуральное число $ p \in \N $ с $ p \geq 2 $, которое делится только на себя и на $1$. Все простые числа, кроме числа $ 2 $, называются \term{нечётными} простыми числами (odd prime numbers). Мы используем $ \Prim $ для обозначения множества всех простых чисел и $ \Prim _{\geq 3} $ для обозначения множества всех нечётных простых чисел.\footnote{Поскольку простые числа имеют центральное значение для нашей темы, заинтересованный читатель может извлечь пользу из широкого спектра книг, посвящённых теории чисел. Введение дано, например, в \chaptname{} 1 и 2 \cite{hardy-2008}. Элементарное школьное введение можно найти в \chaptname{} 33 \cite{wu-1}. Глава 34 \cite{wu-1} даёт введение в основную теорему арифметики (fundamental theorem of arithmetic) \eqref{def:fundamental_theorem_arithmetic}. \cite{fine-2016} может быть особенно интересна для более продвинутых читателей.}

Множество простых чисел $\Prim$ бесконечно, и его можно упорядочить по величине. Это означает, что для любого простого числа $p\in\Prim$ всегда можно найти другое простое число $p'\in\Prim$ с $p<p'$. В результате не существует наибольшего простого числа. Поскольку простые числа можно упорядочить по величине, мы можем записать их следующим образом:
\begin{equation}
\label{eq: primenumber_sequence}
2, 3, 5, 7, 11, 13, 17, 19, 23, 29, 31, 37, 41, 43, 47, 53, 59, 61, 67, \ldots
\end{equation}
Как гласит \term{основная теорема арифметики} (fundamental theorem of arithmetic), простые числа являются в определённом смысле базовыми строительными блоками, из которых составлены все остальные натуральные числа: каждое натуральное число можно получить перемножением простых чисел друг с другом. Чтобы увидеть это, возьмём $ n \in \N $ — любое натуральное число с $n>1$. Тогда всегда существуют простые числа $ p_1, p_2, \ldots, p_k \in \Prim $, такие что выполняется следующее уравнение:
\begin{equation}
\label{def:fundamental_theorem_arithmetic}
n = p_1 \cdot p_2 \cdot \ldots \cdot p_k \;
\end{equation}
Это представление единственно для каждого натурального числа (за исключением порядка \term{множителей} $ p_1, p_2, \ldots, p_k$) и называется \term{разложением на простые множители} (prime factorization) $n$.

\begin{example}[Разложение на простые множители]\label{ex-prime-factorization} Чтобы понять, что мы подразумеваем под разложением на простые множители, рассмотрим число $504\in\N$. Чтобы получить его простые множители, мы можем последовательно делить его на все простые числа в порядке возрастания, начиная с $2$:
\begin{equation*}
504 = 2\cdot 2\cdot 2\cdot 3\cdot 3\cdot 7
\end{equation*}
Мы можем проверить наши результаты, используя Sage, который предоставляет алгоритм разложения натуральных чисел на множители:
\begin{sagecommandline}
sage: n = NN(504)
sage: factor(n)
\end{sagecommandline}
\end{example}

Вычисление из предыдущего примера выявляет важное наблюдение: вычисление разложения на множители целого числа вычислительно затратно, поскольку нам приходится многократно делить на все простые числа, меньшие самого числа, пока все множители не станут простыми. Отсюда возникает важный вопрос: насколько быстро мы можем вычислить разложение на простые множители натурального числа? Этот вопрос представляет собой знаменитую \hilight{проблему разложения целых чисел на множители} (integer factorization problem), и, насколько нам известно, в настоящее время не существует метода, который мог бы разлагать целые числа на множители значительно быстрее наивного подхода, заключающегося в простом делении данного числа на все простые числа в порядке возрастания.

С другой стороны, вычисление произведения заданного множества простых чисел выполняется быстро: нужно просто перемножить все множители. Это простое наблюдение подразумевает, что два процесса — «умножение простых чисел» с одной стороны и обратный процесс «разложения натурального числа на множители» — имеют очень разные вычислительные затраты. Проблема разложения на множители является поэтому примером так называемой \term{односторонней функции} (one-way function): обратимой функции, которую легко вычислить в одном направлении, но сложно в другом
%(see Wikipedia for a description of \href{https://en.wikipedia.org/wiki/Time_complexity}{time complexity} and \href{https://en.wikipedia.org/wiki/Time_complexity#Polynomial_time}{polynomial time})
. 
\footnote{Следует отметить, что то, что «легко» и «сложно» вычислить, зависит от имеющейся в нашем распоряжении вычислительной мощности. Современные компьютеры не могут легко вычислить разложение на простые множители натуральных чисел (в формальных терминах, они не могут вычислить его за полиномиальное время). Однако американский математик Питер В. Шор разработал алгоритм в \citeyear{shor94}, который может вычислить разложение на простые множители натурального числа за полиномиальное время на квантовом компьютере. Последствием этого является то, что криптосистемы, основанные на вычислительной сложности проблемы простых множителей на современных компьютерах, становятся небезопасными, как только становятся доступны практически применимые квантовые компьютеры.}


\begin{exercise}
Какова абсолютная величина целых чисел $-123$, $27$ и $0$?
\end{exercise}
\begin{exercise}
Вычислите разложение на множители $30030$ и проверьте свои результаты с помощью Sage.
\end{exercise}
\begin{exercise}
Рассмотрим следующее уравнение:
\begin{equation*}4\cdot x + 21 = 5\end{equation*}
Вычислите множество всех решений для $x$ при следующих альтернативных предположениях:
\begin{enumerate}
\item Уравнение определено над множеством натуральных чисел.
\item Уравнение определено над множеством целых чисел.
\end{enumerate}
\end{exercise}
\begin{exercise}
Рассмотрим следующее уравнение:
\begin{equation*} 2 x^3 - x^2 - 2 x = - 1.\end{equation*}
Вычислите множество всех решений $x$ при следующих предположениях:
\begin{enumerate}
\item Уравнение определено над множеством натуральных чисел.
\item Уравнение определено над множеством целых чисел.
\item Уравнение определено над множеством рациональных чисел.
\end{enumerate}
\end{exercise}

\subsection{\concept{Деление с остатком}}
\label{Euclidean_division}
Как нам известно из школьной математики, целые числа можно складывать, вычитать и умножать, и результат этих операций гарантированно также будет целым числом.
С другой стороны, деление (в обычном понимании) не определено для целых чисел, поскольку, например, $7$, разделённое на $3$, не даст целого числа. Однако всегда возможно разделить любые два целых числа, если рассматривать \term{деление с остатком} (division with a remainder). Например, $7$, разделённое на $3$, равно $2$ с остатком $1$, поскольку $7 = 2\cdot 3 + 1$.

В этом разделе вводится деление с остатком для целых чисел, обычно называемое \term{\concept{делением с остатком}} (Euclidean division). Это важнейшая техника, лежащая в основе многих понятий этой книги.\footnote{\concept{Деление с остатком} (Euclidean division) вводится в \chaptname{} 1, \secname{} 5 \cite{mignotte-1992} и в \chaptname{} 1, \secname{} 1.3 \cite{cohen-2010}.} Точное определение следующее:

Пусть $ a \in \Z $ и $ b \in \Z $ — два целых числа с $b\neq 0$. Тогда всегда существует другое целое число $ m \in \Z $ и неотрицательное целое число $ r \in \N_0 $ с $ 0 \leq r <|b| $, такие что выполняется следующее:
\begin{equation}
\label{eq_euclidean_division}
a = m \cdot b + r
\end{equation}
Это разложение $a$ при данном $b$ называется \term{\concept{делением с остатком}} (Euclidean division), где $ a $ называется \term{делимым} (dividend), $ b $ — \term{делителем} (divisor), $m$ — \term{частным} (quotient), а $r$ — \term{остатком} (remainder). Можно показать, что и частное, и остаток всегда существуют и единственны, пока делитель отличен от $0$.
\begin{notation}
\label{eq_euclidean_division_notation}
Предположим, что числа $a$, $b$, $m$ и $r$ удовлетворяют уравнению (\ref{eq_euclidean_division}). Тогда мы можем использовать следующие обозначения для частного и остатка \concept{деления с остатком}:
\begin{equation}
\label{def_integer_division_and_modulus}
\begin{array}{lcr}
\Zdiv{a}{b}: = m, & & \Zmod{a}{b}: = r
\end{array}
\end{equation}

Мы также говорим, что целое число $ a $ \term{делится} (divisible) на другое целое число $ b $, если выполняется $ \Zmod{a}{b} = 0 $. В этом случае мы пишем $ b | a $ и называем целое число $\Zdiv{a}{b}$ \term{кофактором} (cofactor) $b$ в $a$.\label{def:cofactor}
\end{notation}
Итак, вкратце, \concept{деление с остатком} — это процесс деления одного целого числа на другое таким образом, который даёт частное и неотрицательный остаток, последний из которых меньше абсолютной величины делителя.

\begin{example}
\label{example:euclidean_division_1}
 Поскольку $ -17 = -5 \cdot 4 + 3 $ является \concept{делением с остатком} $-17$ на $4$, применяя \concept{деление с остатком} и обозначения, определённые в \ref{def_integer_division_and_modulus}, к делимому $-17$ и делителю $4$, мы получаем следующее:
\begin{equation}
\begin{array}{lcr}
\Zdiv{-17}{4} = - 5, & & \Zmod{-17}{4} = 3
\end{array}
\end{equation}
Остаток по определению является неотрицательным числом. В данном случае $4$ не делит $-17$, поскольку остаток не равен нулю. Истинностное значение выражения $4 | -17 $ поэтому равно \textsc{false}. С другой стороны, истинностное значение $4 | 12$ равно \bool{true}, поскольку $4$ делит $12$, так как $ \Zmod{12}{4} = 0 $. Если мы используем Sage для вычисления, мы получаем следующее:
\begin{sagecommandline}
sage: ZZ(-17) // ZZ(4) # Integer quotient
sage: ZZ(-17) % ZZ(4) # remainder
sage: ZZ(4).divides(ZZ(-17)) # self divides other
sage: ZZ(4).divides(ZZ(12))
\end{sagecommandline}
\end{example}
\begin{remark} В \ref{def_integer_division_and_modulus} мы определили обозначения \hilight{$\Zdiv{a}{b}$} и \hilight{$\Zmod{a}{b}$} в терминах \concept{деления с остатком}. Однако следует отметить, что многие языки программирования (такие как Python и Sage) реализуют операторы $(/)$ и $(\%)$ иначе. Программисты должны быть осведомлены об этом, поскольку расхождение между математическим обозначением и реализацией в языках программирования может стать источником тонких ошибок в реализациях криптографических примитивов.

Чтобы привести пример, рассмотрим делимое $-17$ и делитель $-4$. Заметим, что, в отличие от предыдущего \examplename{} \ref{example:euclidean_division_1}, теперь у нас отрицательный делитель. Согласно нашему определению мы имеем следующее:
\begin{equation}\label{euclidean-negative}
\begin{array}{lcr}
\Zdiv{-17}{-4} = 5, & & \Zmod{-17}{-4} = 3
\end{array}
\end{equation}
$ -17 = 5 \cdot (-4) + 3 $ является \concept{делением с остатком} $-17$ на $-4$ (остаток по определению является неотрицательным числом). Однако, используя операторы $(/)$ и $(\%)$ в Sage, мы получаем другой результат: 

\begin{sagecommandline}
sage: ZZ(-17) // ZZ(-4) # Integer quotient
sage: ZZ(-17) % ZZ(-4) # remainder
sage: ZZ(-17).quo_rem(ZZ(-4)) # not Euclidean Division
\end{sagecommandline}
\end{remark}

Методы вычисления \concept{деления с остатком} для целых чисел называются \term{алгоритмами целочисленного деления} (integer division algorithms). Вероятно, наиболее известным алгоритмом является так называемое \term{деление столбиком} (long division), которое большинство из нас, возможно, изучало в школе. Обширное элементарное школьное введение в деление столбиком можно найти в \chaptname{} 7 \cite{wu-1}.

Вкратце, алгоритм деления столбиком перебирает цифры делимого слева направо, вычитая на каждом этапе наибольшую возможную кратную делителя (на уровне цифр). Эти кратные затем становятся цифрами частного, а остаток — первой цифрой делимого.

Поскольку деление столбиком — это стандартный метод, используемый для деления «от руки» многозначных чисел, выраженных в десятичной записи, мы используем его на протяжении всей книги, когда выполняем простые вычисления «от руки», поэтому читатели должны овладеть им. Однако вместо формального определения алгоритма мы приведём несколько примеров, поскольку это, надеемся, сделает процесс более понятным.

\begin{example}[Деление целых чисел столбиком] Чтобы привести пример алгоритма деления целых чисел столбиком, разделим целое число $a=143785$ на число $b=17$. Наша цель, следовательно, состоит в том, чтобы найти решения уравнения \ref{eq_euclidean_division}, то есть нам нужно найти частное $m\in\Z$ и остаток $r \in \N$, такие что $143785 = m\cdot 17 + r$. Используя обозначение, которое в основном используется в англоязычных странах, мы вычисляем следующим образом:\smecomb{SB: I think a more detailed explanation is needed for those unfamiliar with this notation/algorithm. MIRCO: I think we have to make a cut somewhere. Explaining third grade integer division is ot of scope and a requirement to read the book}{Add explanation}
\begin{equation}
\intlongdivision{143785}{17}
\end{equation}
Мы вычислили $m=8457$ и $r=16$, и, действительно, уравнение $143785 = 8457\cdot 17 + 16$ выполняется. Мы можем проверить это, используя Sage:
\begin{sagecommandline}
sage: ZZ(143785).quo_rem(ZZ(17))
sage: ZZ(143785) == ZZ(8457)*ZZ(17) + ZZ(16) # check
\end{sagecommandline}

\end{example}
\begin{exercise}[Деление целых чисел столбиком]
Найдите $m\in\Z$ и $r\in\N$ с $0\leq r< |b|$, такие что $a= m\cdot b +r$ выполняется для следующих пар:
\begin{itemize}\compactlist{}
 \item $(a,b) = (27,5)$ 
 \item $(a,b)=(27,-5)$
 \item $(a,b)=(127,0)$
 \item $(a,b)= (-1687, 11)$
 \item $(a,b)= (0, 7)$
 \end{itemize}
 В каких случаях ваши решения единственны?
\end{exercise}
\begin{exercise}[Алгоритм деления столбиком]
Используя язык программирования по вашему выбору, напишите алгоритм, который вычисляет целочисленное деление столбиком и корректно обрабатывает все крайние случаи.
\end{exercise}
\begin{exercise}[Двоичное представление]
Напишите алгоритм, который вычисляет двоичное представление \ref{def:binary_representation_integer} любого неотрицательного целого числа.
\end{exercise}

\subsection{Расширенный алгоритм Евклида}
Одной из наиболее критически важных частей этой книги является модульная арифметика, определённая в разделе \ref{modular_arithmetic}, и её применение в вычислениях \term{простых полей} (prime fields), определённых в разделе \ref{prime_fields}. Чтобы иметь возможность выполнять вычисления в модульной арифметике, мы должны познакомиться с так называемым \term{расширенным алгоритмом Евклида} (Extended Euclidean Algorithm), используемым для вычисления \term{наибольшего общего делителя} (greatest common divisor, GCD) целых чисел.\footnote{Более глубокое введение в содержание этого раздела можно найти в \chaptname{} 1, \secname{} 1.3 \cite{cohen-2010} и в \chaptname{} 1, \secname{} 8 \cite{mignotte-1992}.}

\term{Наибольший общий делитель} (greatest common divisor) двух ненулевых целых чисел $a$ и $b$ определяется как наибольшее ненулевое натуральное число $d$, такое что $d$ делит и $a$, и $b$, то есть $d|a$, а также $d|b$ истинны. Мы используем обозначение $ gcd (a, b):=d $ для этого числа. Поскольку натуральное число $1$ делит любое другое целое число, $1$ всегда является общим делителем любых двух ненулевых целых чисел, но не обязательно наибольшим.

Общим методом вычисления наибольшего общего делителя является так называемый алгоритм Евклида. Однако, поскольку он нам не понадобится в этой книге, мы введём расширенный алгоритм Евклида, который является методом вычисления наибольшего общего делителя двух натуральных чисел $ a $ и $ b \in \N $, а также двух дополнительных целых чисел $ s, t \in \Z $, таких что выполняется следующее уравнение:
\begin{equation}
\label{eq: erw_Eukl_algo}
gcd (a, b) = s \cdot a + t \cdot b
\end{equation}

Псевдокод в алгоритме \ref{alg_ext_euclid_alg} показывает подробно, как вычислить наибольший общий делитель и числа $s$ и $t$ с помощью расширенного алгоритма Евклида:

\begin{algorithm}\caption{Расширенный алгоритм Евклида}
\label{alg_ext_euclid_alg}
\begin{algorithmic}[0]
\Require $a,b \in \N$ with $a\geq b$
\Procedure{Ext-Euclid}{$a,b$}
\State $r_0\gets a$ and $r_1\gets b$
\State $s_0\gets 1$ and $s_1\gets 0$
\State $t_0\gets 0$ and $t_1\gets 1$
\State $k\gets 2$
\While{$ r_{k-1} \neq 0 $}
\State $ q_k\gets \Zdiv{r_{k-2}}{r_{k-1}} $
\State $ r_{k}\gets \Zmod{r_{k-2}}{r_{k-1}} $
\State $ s_{k}\gets s_{k-2} -q_k \cdot s_{k-1} $
\State $ t_{k}\gets t_{k-2} -q_k \cdot t_{k-1} $
\State $ k \gets k + 1 $
\EndWhile
\State \textbf{return} $gcd(a,b)\gets r_{k-2}$, $s\gets s_{k-2}$ and $ t \gets t_{k-2} $
\EndProcedure
\Ensure $ gcd (a, b) = s \cdot a + t \cdot b $
\end{algorithmic}
\end{algorithm}
Алгоритм достаточно прост, чтобы эффективно использоваться в примерах «от руки». Обычно его записывают в виде таблицы, где строки представляют цикл while, а столбцы — значения массивов $r$, $s$ и $t$ с индексом $k$. Следующий пример демонстрирует простое выполнение.
\begin{example}
\label{example:extended_Euclidean_division_1}
Чтобы проиллюстрировать алгоритм \ref{alg_ext_euclid_alg}, применим его к числам $a=12$ и $b=5$. Поскольку $12,5\in \N$ и $12\geq 5$, все требования выполнены, и мы вычисляем следующим образом:
\begin{center}
  \begin{tabular}{c | c c l}
    k & $ r_k $ & $ s_k $ & $ t_k $ \\\hline
    0 & 12 & 1 & 0 \\
    1 & 5 & 0 & 1 \\
    2 & 2 & 1 & -2 \\
    3 & 1 & -2 & 5 \\
    4 & 0 &  &  \\
    5 & $\cdot$ &&\\
  \end{tabular}
  \end{center}

  
  \begin{comment}
  \begin{tabularx}{1.02\textwidth}{c| r  r  r  r@{}c@{\hspace{0.5cm}}>{\small}X>{\small}X>{\small}X>{\small}X@{}}
    k & $ r_k $ & $ s_k $ & $ t_k $ &$q_k$&&&&& \\\cline{1-5}
    0 & 12 & 1 & 0 &--&& $r_0 \gets a = 12$ & $s_0 \gets 1$ & $t_0 \gets 0$ & \\\\
    1 & 5 & 0 & 1 &--&& $r_1 \gets b = 5$ & $s_1 \gets 0$ & $t_1 \gets 1$ & \\\\
    2 & 2 & 1 & -2 &2 && $r_2 \gets \Zmod{r_0}{r_1} = \newline\Zmod{12}{5} = 2 $  & $s_2 \gets s_0 - q_2\cdot s_1 =\newline 1-2\cdot 0 = 1$ & $t_2 \gets t_0 - q_2\cdot t_1 =\newline 0 - 2\cdot 1 = -2$  &$q_2 \gets \Zdiv{r_0}{r_1} = \newline\Zdiv{12}{5} = 2$\\\\
    3 & 1 & -2 & 5 &2 && $r_3 \gets \Zmod{r_1}{r_2} =\newline \Zmod{5}{2} = 1 $  & $s_3 \gets s_1 - q_3\cdot s_2 =\newline 0-2\cdot 1 = -2$ & $t_3 \gets t_1 - q_3\cdot t_2 =\newline 1 - 2\cdot -2 = 5$   &$q_3 \gets \Zdiv{r_1}{r_2} = \newline\Zdiv{5}{2} = 2$\\\\
    4 & 0 & &  & &  &     $r_4 \gets \Zmod{r_2}{r_3} =\newline \Zmod{2}{1}= 0 $&  &  & \\
  \end{tabularx}
  
  
\sme{SB: I'll streamline the outlook of this table once the discrepancy above is resolved.}
\end{comment}

Из этого мы видим, что наибольший общий делитель $12$ и $5$ равен $ gcd (12, 5) = 1 $ и что уравнение $ 1 = (-2) \cdot 12 + 5 \cdot 5 $ выполняется. Мы также можем использовать Sage, чтобы проверить наши результаты:
\begin{sagecommandline}
sage: ZZ(12).xgcd(ZZ(5)) # (gcd(a,b),s,t)
\end{sagecommandline}
\end{example}
\begin{exercise}[Расширенный алгоритм Евклида]
\label{exercise:extended_Euclidean_division_1}
Найдите целые числа $s,t\in\Z$, такие что $gcd(a,b)= s\cdot a +t\cdot b$ выполняется для следующих пар:
\begin{itemize}\compactlist{}
\item $(a,b) = (45,10)$
\item $(a,b)=(13,11)$
\item $(a,b)=(13,12)$
\end{itemize}
\end{exercise}
\begin{exercise}[На пути к простым полям]
\label{exercise_towards_counting_numbers}
Пусть $n\in \N$ — натуральное число и $p$ — простое число, такие что $n<p$. Чему равен наибольший общий делитель $gcd(p,n)$?
\end{exercise}
\begin{exercise}
Найдите все числа $k\in\N$ с $0\leq k \leq 100$, такие что $gcd(100,k) = 5$.
\end{exercise}
\begin{exercise}
Покажите, что $gcd(n,m) = gcd(n+m,m)$ для всех $n,m\in\N$.
\end{exercise}
\subsection{Взаимно простые целые числа}
\label{def:coprime_integers}
\term{Взаимно простые целые числа} (Coprime integers) — это целые числа, которые не имеют общего простого числа в качестве множителя. Как мы увидим в \secname{} \ref{modular_arithmetic}, взаимно простые целые числа важны для наших целей, поскольку в модульной арифметике вычисления, включающие взаимно простые числа, существенно отличаются от вычислений над невзаимно простыми числами (\defname{} \ref{def:modular_arithmetic_rules}).\footnote{Введение в взаимно простые числа можно найти в \chaptname{} 5, \secname{} 1 \cite{hardy-2008}.}

Наивный способ определить, являются ли два целых числа взаимно простыми, состоял бы в последовательном делении обоих чисел на все простые числа, меньшие этих чисел, чтобы увидеть, имеют ли они общий простой множитель. Однако два целых числа взаимно просты тогда и только тогда, когда их наибольший общий делитель равен $1$. Вычисление $gcd$ является поэтому предпочтительным методом, поскольку оно вычислительно более эффективно.

\begin{example} Рассмотрим снова \examplename{} \ref{example:extended_Euclidean_division_1}. Как мы видели, наибольший общий делитель $12$ и $5$ равен $1$. Это означает, что целые числа $12$ и $5$ взаимно просты, поскольку они не имеют общего делителя, кроме $1$, которое не является простым числом.
\end{example}
\begin{exercise} Рассмотрим снова \exercisename{} \ref{exercise:extended_Euclidean_division_1}. Какие пары $(a,b)$ из этого упражнения являются взаимно простыми?
\end{exercise}
\subsection{Представления целых чисел}
\label{sec:integer-rep}
До сих пор мы представляли целые числа в так называемой \term{десятичной позиционной системе}, которая представляет любое целое число как последовательность элементов из множества десятичных цифр $\{0,1,2,3,4,5,6,7,8,9\}$. Однако существуют и другие представления целых чисел, используемые в информатике и криптографии, которые мы хотим выделить:

Так называемая \term{двоичная позиционная система} (или двоичное представление), представляет каждое целое число как последовательность элементов из множества двоичных цифр (или битов) $\{0,1\}$. Более точно, пусть $n\in\NN$ — неотрицательное целое число в десятичном представлении и $b=b_{k-1}b_{k-2}\ldots b_{0}$ — последовательность \term{битов} $b_j\in\{0,1\}\subset\NN$ для некоторого положительного целого числа $k\in\N$. Тогда $b$ является \term{двоичным представлением} $n$, если выполняется следующее уравнение:

\begin{equation}
\label{def:binary_representation_integer}
n = \sum_{j=0}^{k-1} b_j\cdot 2^j
\end{equation}

Мы дополнительно требуем, что либо $n=0$ и $k=1$ с $b_0=0$, либо $n>0$ и $b_{k-1}=1$. Это исключает ведущие нули.

В этом случае мы пишем $Bits(n):= b_{k-1}b_{k-2}\ldots b_{0}$ для двоичного представления $n$, говорим, что $n$ является $k$-битным числом, и называем $k:= |n|_2$ \term{битовой длиной} $n$.

С этим соглашением двоичное представление любого неотрицательного целого числа единственно. Мы называем $b_0$ \term{младшим битом} (least significant bit) и $b_{k-1}$ \term{старшим битом} (most significant bit) и определяем \term{вес Хэмминга} (Hamming weight) целого числа как количество $1$ в его двоичном представлении.\footnote{Больше о двоичном и общем представлении целых чисел по основанию см., например, в \chaptname{} 1 \cite{mignotte-1992}.}

Ещё одно часто используемое представление — так называемая \term{шестнадцатеричная позиционная система}, которая представляет каждое целое число как последовательность элементов из множества из $16$ цифр, обычно записываемых как $\{0,1,2,3,4,5,6,7,8,9,a,b,c,d,e,f\}$.

Если не указано иное, на протяжении всей книги мы используем десятичную позиционную систему для представления чисел, таких как целые числа или рациональные числа. Однако следует отметить, что поскольку реальные криптографические системы часто имеют дело с большими целыми числами, шестнадцатеричная система является общепринятым выбором в таких обстоятельствах, поскольку шестнадцатеричные представления требуют меньше цифр для представления целого числа, чем десятичное представление.
\begin{sagecommandline}
sage: NN(27713).str(2) # Binary representation
sage: ZZ(27713).str(16) # Hexadecimal representation
\end{sagecommandline}
Если дана позиционная система с $k$ цифрами $\{d_0,d_1,\ldots, d_{k-1}\}$, то представление целого числа $d_{j_h}\cdots d_{j_1} d_{j_0}$ в этой системе может быть преобразовано в 
представление целого числа $n$ в десятичной системе по следующему уравнению:
\begin{equation}
\label{eq:from-k-digits-to-deci}
n = \sum_{i=0}^h j_i\cdot k^{i}
\end{equation}
Существуют аналогичные уравнения для преобразования представления целого числа из любой позиционной системы в любую другую. Десятичная система поэтому не является особенной, просто наиболее распространённой. Чтобы справиться с этой неоднозначностью, многие компьютерные системы принимают префиксы к числу, которые указывают, в какой позиционной системе это число выражено. Общие префиксные обозначения:
 
\begin{center}
\begin{tabular}{rl}
 0x & hexadecimal \\
 0b & binary
\end{tabular}
\end{center} 
\begin{example}
Чтобы понять разницу между представлением целого числа и самим целым числом, рассмотрим выражение $11$ как представление целого числа. Заметим, что без ссылки на позиционную систему это выражение неоднозначно. Можно считать, что это выражение представляет целое число $11$ в десятичном представлении. Однако если $11$ рассматривается как выражение в двоичной системе, то оно относится к целому числу $3$ в десятичном представлении. Более того, когда $11$ рассматривается как выражение в шестнадцатеричной системе, то оно относится к целому числу $17$ в десятичном представлении. Поэтому общепринятой практикой является использование префикса для однозначного указания целого числа в позиционной системе. Мы получаем
\begin{align*}
0x11 & = 17\\
0b11 & = 3
\end{align*}
\end{example}
\begin{example} Чтобы увидеть, как уравнение 
\ref{eq:from-k-digits-to-deci} может преобразовать любое представление в десятичное представление, рассмотрим множество шестнадцатеричных цифр $\{0,1,2,3,4,5,6,7,8,9,a,b,c,d,e,f\}$, которое мы можем записать как 
$$\{0_0,1_1,2_2,3_3,4_4,5_5,6_6,7_7,8_8,9_9,a_{10},b_{11},c_{12},d_{13},e_{14},f_{15}\}$$
Если мы хотим преобразовать целое число $y=0x3f7a$, мы можем записать его как $0x3_3f_{15}7_7a_{10}$ и использовать уравнение \ref{eq:from-k-digits-to-deci} для вычисления его десятичного представления. Поскольку шестнадцатеричная система имеет $16$ цифр и $y$ является $4$-значным числом в шестнадцатеричной системе, мы получаем
\begin{align*}
n = & \sum_{i=0}^4 j_i\cdot 16^{i} \\
  = & 10\cdot 16^0+7\cdot 16^1+15\cdot 16^2+3\cdot 16^3\\
  = & 16250
\end{align*}
\end{example}
\begin{exercise}Рассмотрим \term{восьмеричную} позиционную систему, которая представляет целые числа с помощью 8 цифр, обычно записываемых как $\{0,1,2,3,4,5,6,7\}$. Числа в этой системе характеризуются префиксом $0o$. Запишите числа $0o1354$ и $0o777$ в их десятичном представлении.
\end{exercise}

\section{Модульная арифметика}
\label{modular_arithmetic}
\term{Модульная арифметика} (modular arithmetic) — это система целочисленной арифметики, в которой числа «зацикливаются» при достижении определённого значения, подобно тому как вычисления на часах зацикливаются, когда значение превышает число $12$. Например, если часы показывают $11$ часов, то через $20$ часов будет $7$ часов, а не $31$ час. Число $31$ не имеет смысла на обычных часах, показывающих часы.

Число, при котором происходит зацикливание, называется \term{модулем} (modulus). Модульная арифметика обобщает пример с часами на произвольные модули и изучает уравнения и явления, возникающие в этом новом виде арифметики. Она имеет центральное значение для понимания большинства современных криптографических систем, в значительной степени потому, что модульная арифметика предоставляет вычислительную инфраструктуру для алгебраических типов, имеющих криптографически полезные примеры односторонних функций.\

Хотя модульная арифметика кажется очень отличной от обычной целочисленной арифметики, с которой мы все знакомы, мы призываем вас проработать примеры и убедиться, что, привыкнув к идее, что это новый вид вычислений, она покажется гораздо менее пугающей. Подробное введение в модульную арифметику и её приложения в теории чисел можно найти в \chaptname{} 5 – 8 \cite{hardy-2008}. Элементарное школьное введение в части темы этого раздела можно найти в части 4 \cite{wu-1}.

\subsection{Сравнение}
В дальнейшем пусть $n\in\N$ с $n\geq 2$ — фиксированное натуральное число, которое мы будем называть \term{модулем} (modulus) нашей системы модульной арифметики. Имея такое $n$, мы можем сгруппировать целые числа в классы: два целых числа находятся в одном классе, если их \concept{деление с остатком} (\ref{Euclidean_division}) на $n$ даёт одинаковый остаток. Два числа, находящиеся в одном классе, называются \term{сравнимыми} (congruent).

\begin{example}
Если мы выберем $n=12$, как в нашем примере с часами, то целые числа $-7$, $5$, $17$ и $29$ все сравнимы по модулю $12$, поскольку все они имеют остаток $5$, если мы выполняем \concept{деление с остатком} на них на $12$. Представив аналоговые $12$-часовые часы, начиная с $5$ часов и добавив $12$ часов, мы снова окажемся в $5$ часов, представляя число $17$. Действительно, во многих странах 5:00 дня записывается как 17:00. С другой стороны, когда мы вычитаем $12$ часов, мы снова оказываемся в $5$ часов, представляя число $-7$.
\end{example}
Мы можем формализовать эту интуицию о том, что должно означать сравнение, в надлежащее определение, используя \concept{деление с остатком} (как объяснялось ранее в \ref{integer_arithmetic}). Пусть $ a $, $ b \in \Z $ — два целых числа, и $ n \in \N $ — натуральное число, такое что $n\geq 2$. Целые числа $ a $ и $ b $ называются \term{сравнимыми по модулю} (congruent with respect to the modulus) $ n $ тогда и только тогда, когда выполняется следующее уравнение:
\begin{equation}
\Zmod{a}{n} = \Zmod{b}{n}
\end{equation}

Если, с другой стороны, два числа не сравнимы по данному модулю $n$, мы называем их \term{несравнимыми} (incongruent) относительно $n$.

Иными словами, \term{сравнение} (congruence) — это уравнение «по модулю», что означает, что уравнение должно выполняться только если мы берём модуль обеих частей. Это выражается следующим обозначением:
\footnote{Более глубокое введение в понятие сравнимости, её основные свойства и
приложения в теории чисел можно найти в \chaptname{} 5 \cite{hardy-2008}.}
\begin{equation}
\kongru{a}{b}{n}
\end{equation}
\begin{exercise}
Какие из следующих пар чисел сравнимы по модулю $13$:
\begin{itemize}\compactlist
\item $(5,19)$
\item $(13,0)$
\item $(-4,9)$
\item $(0,0)$
\end{itemize}
\end{exercise}
\begin{exercise}
Найдите все целые числа $x$, такие что сравнение $\kongru{x}{4}{6}$ выполняется.
\end{exercise}
\subsection{Правила вычислений} Определив понятие сравнения как уравнение «по модулю», возникает вопрос: можем ли мы манипулировать сравнением аналогично уравнению? Действительно, мы можем почти применять те же правила подстановки к сравнению, что и к уравнению, с главным отличием, что для некоторого ненулевого целого числа $k\in \Z$ сравнение $\kongru{a}{b}{n}$ эквивалентно сравнению $\kongru{k\cdot a}{k\cdot b}{n}$ только если $k$ и модуль $n$ взаимно просты (см. \ref{def:coprime_integers}). 

Предположим, что даны целые числа $a_1,a_2,b_1,b_2, k\in\Z$ (см. \chaptname{} 5 \cite{hardy-2008}). Тогда для сравнений выполняются следующие арифметические правила:
\begin{itemize}
\label{def:modular_arithmetic_rules}
\item $\kongru{a_1}{b_1}{n}\Leftrightarrow \kongru{a_1+k}{b_1+k}{n}$ (совместимость с переносом)
\item $\kongru{a_1}{b_1}{n}\Rightarrow\kongru{k\cdot a_1}{k\cdot b_1}{n}$ (совместимость с масштабированием)
\item $gcd(k,n)=1 \text{ and } \kongru{k\cdot a_1}{k\cdot b_1}{n}\Rightarrow \kongru{a_1}{b_1}{n}$
\item $\kongru{k\cdot a_1}{k\cdot b_1}{k\cdot n}\Rightarrow\kongru{a_1}{b_1}{n}$
\item $\kongru{a_1}{b_1}{n}\text{ and } \kongru{a_2}{b_2}{n}\Rightarrow \kongru{a_1+a_2}{b_1+b_2}{n}$ (совместимость со сложением)
\item $\kongru{a_1}{b_1}{n}\text{ and } \kongru{a_2}{b_2}{n}\Rightarrow\kongru{a_1\cdot a_2}{b_1\cdot b_2}{n}$ (совместимость с умножением)
\end{itemize}
Другие правила, такие как совместимость с вычитанием, следуют из правил выше. Например, совместимость с вычитанием следует из совместимости с масштабированием при $k=-1$ и совместимости со сложением.

Ещё одно свойство сравнений, не встречающееся в традиционной арифметике целых чисел, — это \term{\concept{малая теорема Ферма}} (Fermat's little theorem). Проще говоря, она утверждает, что в модульной арифметике любое число, возведённое в степень простого числа-модуля, сравнимо с самим числом. Или, более точно, если $ p \in \Prim $ — простое число и $ k \in \mathbb{Z} $ — целое число, то выполняется следующее:
\begin{equation}\label{fermats-little-theorem}
\kongru{k ^ p}{k}{p} 
\end{equation}
Если $k$ взаимно просто с $p$, то мы можем разделить обе части этого сравнения на $k$ и переписать выражение в следующей эквивалентной форме:\footnote{Малая теорема Ферма имеет большое значение в теории чисел. Для подробного доказательства и обширного введения в её следствия см., например, \chaptname{} 6 \cite{hardy-2008}.}
\begin{equation}
\label{eq_fermat_lt_2}
\kongru{k ^{p-1}}{1}{p}
\end{equation}
Приведённый ниже код Sage вычисляет примеры \concept{малой теоремы Ферма} и подчёркивает эффекты от того, является ли показатель $k$ взаимно простым с $p$ (как в случае 137 и 64) или нет (как в случае 1918 и 137):
\begin{sagecommandline}
sage: ZZ(137).gcd(ZZ(64))
sage: ZZ(64)^ ZZ(137) % ZZ(137) == ZZ(64) % ZZ(137)
sage: ZZ(64)^ ZZ(137-1) % ZZ(137) == ZZ(1) % ZZ(137)
sage: ZZ(1918).gcd(ZZ(137))
sage: ZZ(1918)^ ZZ(137) % ZZ(137) == ZZ(1918) % ZZ(137)
sage: ZZ(1918)^ ZZ(137-1) % ZZ(137) == ZZ(1) % ZZ(137)
\end{sagecommandline}

Следующий пример содержит большинство понятий, описанных в этом разделе.

\begin{example}
\label{example_first_congruence}
Чтобы лучше понять сравнения, решим следующее сравнение для $x\in \Z$ в арифметике по модулю $6$:
$$\kongru{7\cdot(2x+21) + 11}{x-102}{6}$$
Поскольку многие правила для сравнений более или менее совпадают с правилами для уравнений, мы можем действовать аналогично тому, как мы бы действовали, если бы имели уравнение для решения. Поскольку обе части сравнения содержат обычные целые числа, мы можем переписать левую часть следующим образом: 

$$7\cdot(2x+21) + 11 = 14x + 147 + 11 = 14x +158$$

Используя это выражение, мы можем переписать сравнение в следующей эквивалентной форме:

$$\kongru{14x +158}{x-102}{6}$$

На следующем шаге мы хотим сместить все вхождения $x$ влево, а все остальные члены — вправо. Поэтому мы применяем правило «совместимости с переносом» дважды. На первом шаге выбираем $k=-x$, а на втором — $k=-158$. \sme{SB: let's separate these two steps in the equivalence below -- separate steps 1 and 2} Поскольку «совместимость с переносом» преобразует сравнение в эквивалентную форму, множество решений не изменится, и мы получаем следующее:
\begin{multline*}
\kongru{14x +158}{x-102}{6} \Leftrightarrow\\
\kongru{14x-x +158-158}{x-x-102-158}{6} \Leftrightarrow \\
\kongru{13x}{-260}{6}
\end{multline*}
Если бы наше сравнение было обычным целочисленным уравнением, мы бы разделили обе части на $13$, получив $x=-20$ как решение. Однако в случае сравнения нам нужно убедиться, что модуль и число, на которое мы хотим разделить, взаимно просты, чтобы гарантировать, что результат деления является выражением, эквивалентным исходному (см. правило \ref{eq_fermat_lt_2}\sme{check reference}). Это означает, что нам нужно найти наибольший общий делитель $gcd(13,6)$. Поскольку $13$ — простое число, а $6$ не кратно $13$, мы знаем, что $gcd(13,6)=1$, то есть эти числа действительно взаимно просты. Мы поэтому вычисляем следующим образом:
$$
\kongru{13x}{-260}{6} \Leftrightarrow \kongru{x}{-20}{6}
$$
Наша задача теперь состоит в том, чтобы найти все целые числа $x$, такие что $x$ сравнимо с $-20$ по модулю $6$. Иными словами, мы должны найти все $x$, такие что выполняется следующее уравнение:
$$
\Zmod{x}{6} = \Zmod{-20}{6}
$$
Поскольку $-4\cdot 6 +4 = -20$, мы знаем, что $ \Zmod{-20}{6} = 4$, и следовательно знаем, что $x=4$ является решением этого сравнения. Однако $22$ — это другое решение, поскольку $ \Zmod{22}{6} = 4$ также. Ещё одно решение — $-20$. На самом деле, существует бесконечно много решений, задаваемых следующим множеством:
$$
\{\ldots, -8,-2, 4,10, 16,\ldots\} = \{4+k\cdot 6 \;|\; k\in \Z\}
$$
Собрав всё вместе, мы показали, что каждое $x$ из множества $\{x=4+k\cdot 6 \;|\; k\in \Z\}$ является решением сравнения $\kongru{7\cdot(2x+21) + 11}{x-102}{6}$. Мы проверим для двух произвольных чисел из этого множества, $x=4$ и $x=4 + 12\cdot 6 = 76$, используя Sage:
\begin{sagecommandline}
sage: (ZZ(7)* (ZZ(2)*ZZ(4) + ZZ(21)) + ZZ(11))  % ZZ(6) == (ZZ(4) - ZZ(102))  % ZZ(6)
sage: (ZZ(7)* (ZZ(2)*ZZ(76) + ZZ(21)) + ZZ(11))  % ZZ(6) == (ZZ(76) - ZZ(102))  % ZZ(6)
\end{sagecommandline}
\end{example}
Если вы ещё не знакомы с модульной арифметикой и чувствуете себя обескураженными тем, насколько сложной она кажется на данном этапе, вам следует помнить две вещи. Во-первых, вычисление сравнений в модульной арифметике на самом деле не сложнее вычислений в более знакомых числовых системах (например, рациональных числах), это просто вопрос привыкания. Во-вторых, как только мы введём идею представлений классов вычетов в \ref{def:remainder_class_representation}, вычисления станут концептуально более чёткими и легче управляемыми.
\begin{exercise}
\label{exercise_congruence_in_F13}
Рассмотрим модуль $13$ и найдём все решения $x\in \Z$ следующего сравнения: $$\kongru{5x+4}{28+2x}{13}$$
\end{exercise}
\begin{exercise}Рассмотрим модуль $23$ и найдём все решения $x\in \Z$ следующего сравнения: $$\kongru{69x}{5}{23}$$
\end{exercise}
\begin{exercise}Рассмотрим модуль $23$ и найдём все решения $x\in \Z$ следующего сравнения: $$\kongru{69x}{46}{23}$$
\end{exercise}
\begin{exercise}
Пусть $a,b,k$ — целые числа, такие что $\kongru{a}{b}{n}$ выполняется. Покажите $\kongru{a^k}{b^k}{n}$.
\end{exercise}
\begin{exercise}
Пусть $a,n$ — целые числа, такие что $a$ и $n$ не взаимно просты. Для каких $b\in\Z$ сравнение $\kongru{a\cdot x}{b}{n}$ имеет решение $x$ и как выглядит множество решений в этом случае?
\end{exercise}
\subsection{Китайская теорема об остатках} Мы видели, как решать сравнения в модульной арифметике. В этом разделе мы рассмотрим, как решать системы сравнений с разными модулями, используя \term{китайскую теорему об остатках} (Chinese Remainder Theorem). Она утверждает, что для любого $ k \in \N $ и взаимно простых натуральных чисел $ n_1, \ldots n_k \in \N $, а также целых чисел $ a_1, \ldots a_k \in \Z $, так называемые \term{одновременные сравнения} (simultaneous congruences) (в \ref{eq_simultaneous_congruence}
 ниже) имеют решение, и все возможные решения этой системы сравнений сравнимы по модулю
произведения $N= n_1 \cdot \ldots \cdot n_k $.\footnote{Это классическая китайская теорема об остатках, известная ещё в древнем Китае. При определённых обстоятельствах теорему можно распространить на невзаимно простые модули $ n_1, \ldots, n_k $, но это выходит за рамки этой книги. Заинтересованные читатели должны обратиться к \chaptname{} 1, \secname{} 1.3.3 \cite{cohen-2010} для введения в теорему и её приложения в вычислительной теории чисел. Доказательство теоремы дано, например, в \chaptname{} 1, \secname{} 10 \cite{mignotte-1992}.} 
\begin{equation}
\label{eq_simultaneous_congruence}
\begin{array}{c}
\kongru{x}{a_1}{n_1} \\
\kongru{x}{a_2}{n_2} \\
\cdots \\
\kongru{x}{a_k}{n_k} \\
\end{array}
\end{equation}

Множество решений вычисляется с помощью \algname{} \ref{chinese-remainder-theorem} ниже. 
\begin{algorithm}\caption{Китайская теорема об остатках}
\label{chinese-remainder-theorem}
\begin{algorithmic}[0]
\Require , $k\in \Z$, $j\in \NN$ and $n_0,\ldots,n_{k-1} \in \N$ coprime
\Procedure{Congruence-Systems-Solver}{$a_{0},\ldots,a_{k-1}$}
\State $N\gets n_0\cdot \ldots \cdot n_{k-1}$
\While{$j< k $}
\State $N_j\gets N/n_j$
\State $(\_,s_j,t_j)\gets EXT-EUCLID (N_j,n_j)$
  \Comment{$1 = s_j\cdot N_j + t_j\cdot n_j$}
\EndWhile
\State $x'\gets \sum_{j=0}^{k-1}a_j\cdot s_j\cdot N_j$
\State $x\gets \Zmod{x'}{N}$
\State \textbf{return} $\{x+ m\cdot N\;|\; m\in \Z\}$
\EndProcedure
\Ensure $\{x+ m\cdot N\;|\; m\in \Z\}$ is the complete solution set to \ref{eq_simultaneous_congruence}.
\end{algorithmic}
\end{algorithm}


\begin{example} Чтобы проиллюстрировать, как решать одновременные сравнения с помощью китайской теоремы об остатках, рассмотрим следующую систему сравнений:
$$
\begin{array}{c}
\kongru{x}{4}{7} \\
\kongru{x}{1}{3} \\
\kongru{x}{3}{5} \\
\kongru{x}{0}{11} \\
\end{array}
$$
Очевидно, все модули взаимно просты (поскольку все они простые числа).  Теперь мы вычисляем следующим образом: 
$$
\begin{array}{l}
N = 7 \cdot 3 \cdot 5 \cdot 11 = 1155\\
N_1 = 1155/7 = 165\\
N_2 = 1155/3 = 385\\
N_3 = 1155/5 = 231\\
N_4 = 1155/11 = 105\\
\end{array}
$$

Из этого мы вычисляем с помощью расширенного алгоритма Евклида:\sme{add more explanation}
$$
\begin{array}{cccc}
 1 = & 2 \cdot 165  & + & -47 \cdot 7 \\
 1 = & 1 \cdot 385  & + &  -128 \cdot 3 \\
 1 = & 1 \cdot 231  & + &  -46 \cdot 5 \\
 1 = & 2 \cdot 105  & + &  -19 \cdot 11 \\
\end{array}
$$
В результате мы получаем $x = 4 \cdot 2 \cdot 165 + 1 \cdot 1 \cdot 385 + 3 \cdot 1 \cdot 231 + 0 \cdot 2 \cdot 105 = 2398$
как одно решение. Поскольку $ \Zmod{2398}{1155} = 88, $ множество всех решений равно
$ \{\ldots, -2222, -1067,88,1243, 2398, \ldots \} $. Мы можем использовать вычисление Sage китайской теоремы об остатках (CRT), чтобы проверить наши результаты:
\begin{sagecommandline}
sage: CRT_list([4,1,3,0], [7,3,5,11])
\end{sagecommandline}
\end{example}

\subsection{Представление классом вычетов}
\label{def:remainder_class_representation}
Как мы видели в различных примерах ранее, вычисление сравнений может быть громоздким, а множества решений в общем случае велики. Поэтому выгодно найти некоторое упрощение для модульной арифметики.

К счастью, это возможно и относительно просто, как только мы отождествим каждое множество чисел, имеющих равные остатки, с самим этим остатком, и назовём это множество \term{классом вычетов} (remainder class) или \term{классом вычетов} (residue class) в арифметике по модулю $n$.

Тогда из свойств \concept{деления с остатком} следует, что существует ровно $ n $ различных классов вычетов для каждого модуля $n$, и что целочисленное сложение и умножение могут быть спроецированы на новый вид сложения и умножения на этих классах.

Неформально говоря, новые правила сложения и умножения затем вычисляются путём взятия любого элемента первого класса вычетов и некоторого элемента второго класса вычетов, затем их сложения или умножения обычным способом и определения, в каком классе вычетов содержится результат. Следующий пример делает это абстрактное описание более конкретным.

\begin{example} [Арифметика по модулю $6$]
\label{def_residue_ring_z_6}
Выбрав модуль $ n = 6 $, мы имеем шесть классов вычетов целых чисел, которые сравнимы по модулю $ 6 $, то есть они имеют одинаковый остаток при делении на $6$. Когда мы отождествляем каждый из этих классов вычетов с остатком, мы получаем следующее отождествление:
$$
\begin{array}{l}
0: = \{\ldots, -6,0,6,12, \ldots \}\\
1: = \{\ldots, -5,1,7,13, \ldots \}\\
2: = \{\ldots, -4,2,8,14, \ldots \} \\
3: = \{\ldots, -3,3,9,15, \ldots \}\\
4: = \{\ldots, -2,4,10,16, \ldots \}\\
5: = \{\ldots, -1,5,11,17, \ldots \}
\end{array}
$$
Чтобы вычислить новый закон сложения этих представителей классов вычетов, скажем $2+5$, мы выбираем произвольный элемент из каждого класса, скажем $14$ и $-1$, складываем эти числа обычным способом, а затем смотрим на класс вычетов результата.

Складывая $14$ и $(-1)$, мы получаем $13$, а $13$ находится в классе вычетов $1$. Следовательно, мы находим, что $2+5=1$ в арифметике по модулю $6$, что является более читаемым способом записи сравнения $\kongru{2+5}{1}{6}$.

Применяя то же рассуждение ко всем классам вычетов, сложение и умножение могут быть перенесены на представителей классов вычетов. Результаты для арифметики по модулю $6$ сведены в следующие таблицы сложения (addition tables) и умножения (multiplication tables):
\begin{equation}\label{eq:Z6_tables}
  \begin{tabular}{c | c c c c c c}
    + & 0 & 1 & 2 & 3 & 4 & 5\\\hline
    0 & 0 & 1 & 2 & 3 & 4 & 5 \\
    1 & 1 & 2 & 3 & 4 & 5 & 0\\
    2 & 2 & 3 & 4 & 5 & 0 & 1\\
    3 & 3 & 4 & 5 & 0 & 1 & 2\\
    4 & 4 & 5 & 0 & 1 & 2 & 3\\
    5 & 5 & 0 & 1 & 2 & 3 & 4
  \end{tabular} \quad \quad \quad \quad
  \begin{tabular}{c | c c c c c c}
$ \cdot $ & 0 & 1 & 2 & 3 & 4 & 5 \\\hline
        0 & 0 & 0 & 0 & 0 & 0 & 0\\
        1 & 0 & 1 & 2 & 3 & 4 & 5\\
        2 & 0 & 2 & 4 & 0 & 2 & 4\\
        3 & 0 & 3 & 0 & 3 & 0 & 3\\
        4 & 0 & 4 & 2 & 0 & 4 & 2\\
        5 & 0 & 5 & 4 & 3 & 2 & 1
  \end{tabular}
\end{equation}
Таким образом, мы определили новую арифметическую систему, которая содержит всего $6$ чисел и снабжена собственным определением сложения и умножения. Мы называем её \term{арифметикой по модулю 6} (modular 6 arithmetic) и записываем соответствующий тип как $\Z_6$.

Чтобы увидеть, почему отождествление класса вычетов с его остатком полезно и на самом деле значительно упрощает вычисления сравнений, вернёмся к сравнению из \examplename{} \ref{example_first_congruence}:
\begin{equation}
\label{example_congruence_two_1}
\kongru{7\cdot(2x+21) + 11}{x-102}{6}
\end{equation}

Как показано в \examplename{} \ref{example_first_congruence}, арифметика сравнений может отклоняться от обычной арифметики: например, при делении нужно проверять, являются ли модуль и делимое взаимно простыми, а решения в общем случае не единственны.

Мы можем переписать сравнение в \eqref{example_congruence_two_1} как \hilight{уравнение} над нашим новым арифметическим типом $\Z_6$, \term{проецируя на классы вычетов} (projecting onto the remainder classes): поскольку $\Zmod{7}{6}= 1$, $\Zmod{21}{6}= 3$, $\Zmod{11}{6}= 5$ и $\Zmod{102}{6}= 0$, мы получаем следующее:
\begin{align*}
\kongru{7\cdot(2x+21) + 11}{x-102}{6} \text{ over } \Z\\
\Leftrightarrow & \; \; 1\cdot (2x+3) + 5 = x \text{ over } \Z_6
\end{align*}
Мы можем использовать таблицу умножения и сложения в \eqref{eq:Z6_tables} выше, чтобы решить уравнение справа, как мы бы решали обычные целочисленные уравнения:

\begin{align*}
1\cdot (2x+3) + 5 &= x & \text{ }\\
2x+3 + 5 &= x & \text{\# таблица сложения: } 3+5 = 2 \\
2x+2 &= x & \text{\# прибавим 4 и $-x$ к обеим сторонам} \\
2x+2 +4 -x &= x + 4 -x & \text{\# таблица сложения: } 2+4 = 0 \\
x &= 4 &
\end{align*}

Как мы видим, несмотря на несколько непривычные правила сложения и умножения, решение сравнений таким способом очень похоже на решение обычных уравнений. И, действительно, множество решений совпадает с множеством решений исходного сравнения, поскольку $4$ отождествляется с множеством $\{4+6\cdot k\;|\; k\in\Z\}$.

Мы можем использовать Sage для вычислений в нашем арифметическом типе по модулю $6$. Это особенно полезно для проверки наших вычислений:
\begin{sagecommandline}
sage: Z6 = Integers(6)
sage: Z6(2) + Z6(5)
sage: Z6(7)*(Z6(2)*Z6(4)+Z6(21))+Z6(11) == Z6(4) - Z6(102)
\end{sagecommandline}
\end{example}

\begin{remark}[$k$-битный модуль] В криптографических работах мы иногда читаем фразы вроде «$[\ldots]$ с использованием $4096$-битного модуля». Это означает, что базовый модуль $n$ модульной арифметики, используемой в системе, имеет двоичное представление длиной $4096$ бит. Напротив, число $6$ имеет двоичное представление $110$, и следовательно наш \examplename{} \ref{def_residue_ring_z_6}
 описывает $3$-битную арифметическую систему с модулем.
\end{remark}
\begin{exercise}
Определите $\Z_{13}$ как арифметику по модулю $13$\sme{modulo/ modulus/ modular? unify throughout} аналогично \examplename{} \ref{def_residue_ring_z_6}. Затем рассмотрите сравнение из \exercisename{} \ref{exercise_congruence_in_F13} и перепишите его в виде уравнения в $\Z_{13}$.
\end{exercise}

\subsection{Модульные обратные элементы}
\label{sec:modular_inverses}
Как нам известно, целые числа можно складывать, вычитать и умножать так, что результат также является целым числом, но это не верно для деления целых чисел в общем случае: например, $3/2$ — не целое число. Чтобы понять, почему это так, с более теоретической точки зрения, рассмотрим сначала определение мультипликативно обратного элемента. Когда у нас есть множество, на котором определено некоторое умножение, и у нас есть выделенный элемент этого множества, который ведёт себя нейтрально относительно этого умножения (не изменяет ничего при умножении на любой другой элемент), то мы можем определить \term{мультипликативно обратные элементы} (multiplicative inverses) следующим образом:

\begin{definition}

Пусть $S$ — наше множество, на котором определено некоторое понятие умножения $a\cdot b$ и \term{нейтральный элемент} (neutral element) $1\in S$, такой что $1\cdot a = a$ для всех элементов $a\in S$. Тогда \term{мультипликативно обратный элемент} (multiplicative inverse) $a^{-1}$ элемента $a\in S$ определяется следующим образом:
\begin{equation}
a\cdot a^{-1} = 1
\end{equation}
\end{definition}
Неформально говоря, определение мультипликативно обратного элемента означает, что он «уничтожает» исходный элемент, так что умножение двух даёт $1$.

Числа, имеющие мультипликативно обратные элементы, представляют особый интерес, поскольку они непосредственно приводят к определению деления на эти числа. Действительно, если $a$ — число, для которого существует мультипликативно обратный элемент $a^{-1}$, то мы определяем \term{деление} (division) на $a$ просто как умножение на обратный элемент:
\begin{equation}
\frac{b}{a}:= b\cdot a^{-1}
\end{equation}
\begin{example} Рассмотрим множество рациональных чисел, также известных как дроби, $\mathbb{Q}$. Для этого множества нейтральным элементом умножения является $1$, поскольку $1\cdot a = a$ для всех рациональных чисел. Например, $1\cdot 4=4$, $1\cdot \frac{1}{4}=\frac{1}{4}$, или $1\cdot 0 =0$ и так далее.

Каждое рациональное число $a\neq 0$ имеет мультипликативно обратный элемент, равный $\frac{1}{a}$.
Например, мультипликативно обратный элемент $3$ равен $\frac{1}{3}$, поскольку $3\cdot \frac{1}{3}=1$, мультипликативно обратный элемент $\frac{5}{7}$ равен $\frac{7}{5}$, поскольку $\frac{5}{7}\cdot \frac{7}{5}=1$, и так далее.
\end{example}
\begin{example} Рассмотрим множество $\Z$ целых чисел. Мы видим, что нейтральным элементом умножения является число $1$. Мы также видим, что никакое целое число, кроме $1$ или $-1$, не имеет мультипликативно обратного элемента, поскольку уравнение $a\cdot x =1$ не имеет целочисленных решений для $a\neq 1$ или $a\neq -1$.

Определение мультипликативно обратного элемента имеет аналогичное определение для сложения, называемое \term{аддитивно обратным элементом} (additive inverse). В случае целых чисел нейтральным элементом относительно сложения является $0$, поскольку $a+0=a$ для всех целых чисел $a\in\Z$. Аддитивно обратный элемент всегда существует и равен отрицательному числу $-a$, поскольку $a+(-a)=0$.
\end{example}
\begin{example} Рассмотрим множество $\Z_6$ классов вычетов по модулю $6$ из \examplename{} \ref{def_residue_ring_z_6}. Мы можем использовать таблицу умножения в \eqref{eq:Z6_tables}, чтобы найти мультипликативно обратные элементы. Для этого мы смотрим на строку элемента и находим элемент, равный $1$. Если такой элемент существует, элемент этого столбца является мультипликативно обратным. Если, с другой стороны, в строке нет элемента, равного $1$, мы знаем, что элемент не имеет мультипликативно обратного.

Например, в $\Z_6$ мультипликативно обратный элемент $5$ равен самому $5$, поскольку $5\cdot 5=1$. Мы также видим, что $5$ и $1$ — единственные элементы, имеющие мультипликативно обратные элементы в $\Z_6$.

Теперь, поскольку $5$ имеет мультипликативно обратный элемент в арифметике по модулю $6$, мы можем делить на $5$ в $\Z_6$, поскольку у нас есть понятие мультипликативно обратного элемента, а деление — не что иное, как умножение на мультипликативно обратный элемент:
$$
\frac{4}{5}= 4\cdot 5^{-1} = 4\cdot 5 = 2
$$
\end{example}
Из последнего примера мы можем сделать интересное наблюдение: хотя $5$ не имеет мультипликативно обратного элемента как целое число, оно имеет мультипликативно обратный элемент в арифметике по модулю $6$.

Это поднимает вопрос о том, какие числа имеют мультипликативно обратные элементы в модульной арифметике. Ответ состоит в том, что в арифметике по модулю $n$ число $r$ имеет мультипликативно обратный элемент тогда и только тогда, когда $n$ и $r$ взаимно просты. Поскольку в этом случае $gcd(n,r)=1$, мы знаем из расширенного алгоритма Евклида, что существуют числа $s$ и $t$, такие что выполняется следующее уравнение:
\begin{equation}
\label{eq_compute_multiplicative_inverse}
1 = s\cdot n + t\cdot r
\end{equation}
Если мы берём модуль $n$ обеих частей, член $s\cdot n$ исчезает, что говорит нам о том, что $\Zmod{t}{n}$ является мультипликативно обратным элементом $r$ в арифметике по модулю $n$.\sme{expand on this}
\begin{example}[Мультипликативно обратные элементы в $\Z_6$] В предыдущем примере мы находили мультипликативно обратные элементы в $\Z_6$ по справочной таблице в \eqref{eq:Z6_tables}. В реальных примерах обычно невозможно записать эти справочные таблицы, поскольку модуль слишком велик, а множества иногда содержат больше элементов, чем атомов в наблюдаемой вселенной.

Теперь, пытаясь определить, что $2\in \Z_6$ не имеет мультипликативно обратного элемента в $\Z_6$ без использования справочной таблицы, мы сразу замечаем, что $2$ и $6$ не взаимно просты, поскольку их наибольший общий делитель равен $2$. Отсюда следует, что уравнение \ref{eq_compute_multiplicative_inverse} не имеет решений $s$ и $t$, что означает, что $2$ не имеет мультипликативно обратного элемента в $Z_6$.

То же рассуждение работает для $3$ и $4$, поскольку ни одно из них не взаимно просто с $6$. Случай $5$ отличается, поскольку $gcd(6,5)=1$. Чтобы вычислить мультипликативно обратный элемент $5$, мы используем расширенный алгоритм Евклида и вычисляем следующее:
\begin{center}
  \begin{tabular}{c | c c l}
    k & $ r_k $ & $ s_k $ & $ t_k = \Zdiv{(r_k-s_k \cdot a)}{b} $ \\\hline
    0 & 6 & 1 & 0 \\
    1 & 5 & 0 & 1 \\
    2 & 1 & 1 & -1 \\
    3 & 0 & . & . \\
  \end{tabular}
\end{center}

Мы получаем $s=1$, а также $t=-1$, и имеем $1 = 1\cdot 6 -1\cdot 5$. Из этого следует, что $\Zmod{-1}{6}=5$ является мультипликативно обратным элементом $5$ в арифметике по модулю $6$. Мы можем проверить это с помощью Sage:
\begin{sagecommandline}
sage: ZZ(6).xgcd(ZZ(5))
\end{sagecommandline}
\end{example}
На этом этапе внимательный читатель может заметить, что использование модуля, являющегося простым числом, представляет особый интерес. Как нам известно из \exercisename{} \ref{exercise_towards_counting_numbers}, если $n$ — простое число, то $gcd(r,n)=1$ для каждого $r<n$. Следовательно, все ненулевые классы вычетов по модулю $n$ имеют мультипликативно обратные элементы в этой постановке. Действительно, \concept{малая теорема Ферма} \eqref{fermats-little-theorem} предоставляет способ вычисления мультипликативно обратных элементов в этой ситуации, поскольку в случае простого модуля $p$ и $r<p$ мы получаем следующее:
\begin{align*}
\kongru{r^p}{r}{p} & \Leftrightarrow \\
\kongru{r^{p-1}}{1}{p} & \Leftrightarrow \\
\kongru{r\cdot r^{p-2}}{1}{p}
\end{align*}
Это говорит нам, что мультипликативно обратный элемент класса вычетов $r$ в арифметике по модулю $p$ в точности равен $r^{p-2}$.
\begin{example} [Арифметика по модулю $5$]
\label{primfield_z_5}
Чтобы увидеть уникальные свойства модульной арифметики, когда модуль является простым числом, мы воспроизведём наши результаты из \examplename{} \ref{def_residue_ring_z_6}, но на этот раз для простого модуля $5$.
Для $ p = 5 $ мы имеем пять классов эквивалентности целых чисел, которые сравнимы по модулю $ 5 $. Мы записываем это следующим образом:
$$
\begin{array}{ccc}
0: = \{\ldots, -5,0,5,10, \ldots \}\\
1: = \{\ldots, -4,1,6,11, \ldots \}\\
2: = \{\ldots, -3,2,7,12, \ldots \} \\
3: = \{\ldots, -2,3,8,13, \ldots \}\\
4: = \{\ldots, -1,4,9,14, \ldots \}
\end{array}
$$
Сложение и умножение могут быть перенесены на классы эквивалентности способом, в точности параллельным \examplename{} \ref{def_residue_ring_z_6}. Это даёт следующие таблицы сложения и умножения:
\begin{equation}\label{Z5_tables}
  \begin{tabular}{c | c c c c c}
    + & 0 & 1 & 2 & 3 & 4 \\\hline
    0 & 0 & 1 & 2 & 3 & 4 \\
    1 & 1 & 2 & 3 & 4 & 0 \\
    2 & 2 & 3 & 4 & 0 & 1 \\
    3 & 3 & 4 & 0 & 1 & 2 \\
    4 & 4 & 0 & 1 & 2 & 3 \\
  \end{tabular} \quad \quad \quad \quad
  \begin{tabular}{c | c c c c c}
$ \cdot $ & 0 & 1 & 2 & 3 & 4 \\\hline
      0 & 0 & 0 & 0 & 0 & 0 \\
      1 & 0 & 1 & 2 & 3 & 4 \\
      2 & 0 & 2 & 4 & 1 & 3 \\
      3 & 0 & 3 & 1 & 4 & 2 \\
      4 & 0 & 4 & 3 & 2 & 1 \\
  \end{tabular}
\end{equation}
Называя множество классов вычетов в арифметике по модулю $5$ с этим сложением и умножением $\Z_5$, мы видим некоторые тонкие, но важные отличия от ситуации в $\Z_6$. В частности, мы видим, что в таблице умножения каждый остаток $r\neq 0$ имеет элемент $1$ в своей строке и следовательно имеет мультипликативно обратный элемент. Кроме того, не существует ненулевых элементов, произведение которых равно нулю.

Чтобы использовать \concept{малую теорему Ферма} в $\Z_5$ для вычисления мультипликативно обратных элементов (вместо использования таблицы умножения), рассмотрим $3\in\Z_5$. Мы знаем, что мультипликативно обратный элемент даётся классом вычетов, содержащим $3^{5-2}=3^3=3\cdot 3\cdot 3= 4\cdot 3 = 2$. И действительно $3^{-1}=2$, поскольку $3\cdot 2 =1$ в $\Z_5$.

Мы можем использовать Sage для вычислений в нашем арифметическом типе по модулю $5$, чтобы проверить наши вычисления:
\begin{sagecommandline}
sage: Z5 = Integers(5)
sage: Z5(3)**(5-2)
sage: Z5(3)**(-1)
sage: Z5(3)**(5-2) == Z5(3)**(-1)
\end{sagecommandline}
\end{example}
\begin{example}
Чтобы понять одно из главных отличий между модульной арифметикой с простым модулем и модульной арифметикой с непростым модулем, рассмотрим линейное уравнение $a\cdot x +b=0$, определённое над обоими типами $\Z_5$ и $\Z_6$. Поскольку каждый ненулевой элемент имеет мультипликативно обратный элемент в $\Z_5$, мы всегда можем решать эти уравнения в $\Z_5$, что неверно в $\Z_6$. Чтобы увидеть это, рассмотрим уравнение $3x+3=0$. В $\Z_5$ мы имеем следующее:
\begin{align*}
3x+3    &= 0 & \text{\# прибавим 2 к обеим сторонам} \\
3x+3+2  &= 2 & \text{\# таблица сложения: } 2+3 = 0 \\
3x      &= 2 & \text{\# разделим на } 3  \text{ (что равно умножению на 2)}\\
2\cdot(3x)      &= 2\cdot 2 & \text{\# таблица умножения: } 2\cdot 2=4 \\
 x      &= 4 &
\end{align*}
Так что в случае нашей модульной арифметики с простым модулем мы получаем единственное решение $x=4$. Теперь рассмотрим $\Z_6$:
\begin{align*}
3x+3    &= 0 & \text{\# прибавим 3 к обеим сторонам} \\
3x+3+3  &= 3 & \text{\# таблица сложения: } 3+3 = 0 \\
3x      &= 3 & \text{\# деление невозможно (не существует мультипликативно обратного элемента 3)}
\end{align*}
Так что в этом случае мы не можем решить уравнение относительно $x$, разделив на $3$. И, действительно, когда мы смотрим на таблицу умножения $\Z_6$ (\examplename{} \ref{def_residue_ring_z_6}), мы находим, что существует три решения $x\in\{1,3,5\}$, таких что $3x+3=0$ выполняется для всех них.
\end{example}
\begin{exercise}
Рассмотрим модуль $n=24$. Какие из целых чисел $7$, $1$, $0$, $805$, $-4255$ имеют мультипликативно обратные элементы в арифметике по модулю $24$? Вычислите обратные элементы, если они существуют.
\end{exercise}
\begin{exercise}
Найдите множество всех решений сравнения $\kongru{17(2x+5)-4}{2x+4}{5}$. Затем спроецируйте сравнение в $\Z_5$ и решите полученное уравнение в $\Z_5$. Сравните результаты.
\end{exercise}
\begin{exercise}
Найдите множество всех решений сравнения $\kongru{17(2x+5)-4}{2x+4}{6}$. Затем спроецируйте сравнение в $\Z_6$ и попытайтесь решить полученное уравнение в $\Z_6$.
\end{exercise}
\section{Арифметика многочленов}
\label{sec:polynomial_arithmetics}
Многочлен — это выражение, состоящее из переменных (также называемых неопределёнными) и коэффициентов, которое включает только операции сложения, вычитания и умножения. Все коэффициенты многочлена должны иметь один и тот же тип, например, они должны быть все целыми числами или все рациональными числами и т.д.\footnote{Введение в теорию многочленов можно найти, например, в \chaptname{} 3 \cite{mignotte-1992}, а подробное описание многих алгоритмов, используемых в вычислениях с многочленами, дано в \chaptname{} 3 \cite{cohen-2010}.}

Более точно, \term{одночленный многочлен}\footnote{В нашем контексте термин одночленный означает, что многочлен содержит только одну переменную.} — это выражение, показанное ниже:
\begin{equation}\label{eq:polynomial}
P(x) := \sum _{j = 0} ^{m}{a} _{j}{x} ^{j} ={a} _{m}x^m +{a} _{m-1} x^{m-1} + \dots + a_1 x + a_0 \;,
\end{equation}

Если $P$ не является нулевым многочленом, мы требуем $a_m\neq 0$ и берём $m$ как наибольший индекс с ненулевым коэффициентом. В \eqref{eq:polynomial} $x$ называется \term{переменной} (variable), а каждый $ a$ называется \term{коэффициентом} (coefficient). Если $R$ — тип коэффициентов, то множество всех \hilight{одночленных многочленов с коэффициентами в $R$} записывается как $R[x]$. Одночленные многочлены часто просто называются многочленами и записываются как $ P (x) \in R[x]$. Свободный член $a_0$ также записывается как $ P(0)$.

Многочлен называется \term{нулевым многочленом} (zero polynomial), если все его коэффициенты равны нулю. Многочлен называется \term{единичным многочленом} (one polynomial), если свободный член равен $1$, а все остальные коэффициенты равны нулю.

Дан одночленный многочлен $P(x)=\sum_{j=0}^m a_jx^j$, не являющийся нулевым многочленом, мы называем неотрицательное целое число $deg(P):=m$ \textit{степенью} (degree) $P$, и определяем степень нулевого многочлена как $-\infty$, где $-\infty$ (отрицательная бесконечность) — символ со свойствами $-\infty + m = -\infty$ и $-\infty < m$ для всех неотрицательных целых чисел $m\in\NN$. 

Кроме того, мы обозначаем коэффициент члена с наивысшей степенью, называемый \term{старшим коэффициентом} (leading coefficient), многочлена $P$ следующим образом:\sme{add explanation on why this is important}
\begin{equation}
\label{def_leading_coefficient}
Lc(P):=a_m
\end{equation}

Мы можем ограничить множество $R[x]$ \hilight{всех} многочленов с коэффициентами в $R$ множеством всех таких многочленов, степень которых не превышает определённого значения. Если $m$ — максимально допустимая степень, мы записываем $R_{\leq m}[x]$ для множества всех многочленов со степенью, меньшей или равной $m$.

\begin{example}[Многочлены с целыми коэффициентами]
\label{example:integer_polynomials} Коэффициенты многочлена должны все иметь один и тот же тип. Множество многочленов с целыми коэффициентами записывается как $\Z[x]$. Ниже приведены некоторые примеры таких многочленов:
\begin{align*}
P_1(x) &= 2x^2 -4x +17 & \text{ \# со } deg(P_1)=2 \text{ и } Lc(P_1)=2\\
P_2(x) &= x^{23} & \text{ \# со } deg(P_2)=23 \text{ и } Lc(P_2)=1\\
P_3(x) &= x & \text{ \# со }  deg(P_3)=1 \text{ и } Lc(P_3)=1\\
P_4(x) &= 174 & \text{ \# со }  deg(P_4)=0 \text{ и } Lc(P_4)=174\\
P_5(x) &= 1 & \text{ \# со }  deg(P_5)=0 \text{ и } Lc(P_5)=1\\
P_6(x) &= 0 & \text{ \# со }  deg(P_6)=-\infty \text{ и } Lc(P_6)=0\\
P_7(x) &= (x-2)(x+3)(x-5)
\end{align*}

Каждое целое число можно рассматривать как многочлен с целыми коэффициентами нулевой степени. $P_7$ — многочлен, поскольку мы можем раскрыть его определение в $P_7(x)=x^3 - 4 x^2 - 11 x + 30$, что является многочленом степени $3$ со старшим коэффициентом $1$. 

Следующие выражения не являются многочленами с целыми коэффициентами:
\begin{align*}
Q_1(x) &= 2x^2 + 4 + 3x^{-2}\\
Q_2(x) &= 0.5x^4 -2x\\
Q_3(x) &=2^x
\end{align*}

$Q_1$ не является многочленом с целыми коэффициентами, поскольку выражение $x^{-2}$ имеет отрицательный показатель. $Q_2$ не является многочленом с целыми коэффициентами, поскольку коэффициент $0.5$ не является целым числом. $Q_3$ не является многочленом с целыми коэффициентами, поскольку переменная появляется в показателе степени коэффициента.

Мы можем использовать Sage для вычислений с многочленами. Для этого мы должны указать символ для переменной и тип для коэффициентов. \smecomb{(For the definition of rings see \ref{sec:rings}.)}{why is this ref here?} \smecomb{Note, however, that  Sage defines the degree of the zero polynomial to be $-1$.}{what does this imply?}
\begin{sagecommandline}
sage: Zx = ZZ['x'] # integer polynomials with indeterminate x
sage: Zt.<t> = ZZ[] # integer polynomials with indeterminate t
sage: Zx
sage: Zt
sage: p1 = Zx([17,-4,2])
sage: p1
sage: p1.degree()
sage: p1.leading_coefficient()
sage: p2 = Zt(t^23)
sage: p2
sage: p6 = Zx([0])
sage: p6.degree()
\end{sagecommandline}
\end{example}
\begin{example}[Многочлены над $\mathbb{Z}_6$]
\label{example:integer_mod_6_polynomials} Вспомним определение арифметики по модулю $6$ $\Z_6$ из \examplename{} \ref{def_residue_ring_z_6}. Множество всех многочленов с неопределённой $x$ и коэффициентами в $\Z_6$ обозначается как $\Z_6[x]$. Ниже приведены некоторые примеры многочленов из $\Z_6[x]$:
\begin{align*}
P_1(x) &= 2x^2 -4x +5 & \text{ \# со } deg(P_1)=2 \text{ и } Lc(P_1)=2\\
P_2(x) &= x^{23} & \text{ \# со } deg(P_2)=23 \text{ и } Lc(P_2)=1\\
P_3(x) &= x & \text{ \# со }  deg(P_3)=1 \text{ и } Lc(P_3)=1\\
P_4(x) &= 3 & \text{ \# со }  deg(P_4)=0 \text{ и } Lc(P_4)=3\\
P_5(x) &= 1 & \text{ \# со }  deg(P_5)=0 \text{ и } Lc(P_5)=1\\
P_6(x) &= 0 & \text{ \# со }  deg(P_5)=-\infty \text{ и } Lc(P_6)=0\\
P_7(x) &= (x-2)(x+3)(x-5)
\end{align*}
Так же как в предыдущем примере, $P_7$ — многочлен. Однако, поскольку теперь мы работаем с коэффициентами из $\Z_6$, раскрытие $P_7$ вычисляется иначе, поскольку мы должны использовать сложение и умножение в $\Z_6$, как определено в \eqref{eq:Z6_tables}\sme{check reference}. Мы получаем следующее:
\begin{align*}
(x-2)(x+3)(x-5) &= (x+4)(x+3)(x+1) & \text{\# аддитивно обратные элементы в } \Z_6 \\
                &= (x^2+4x+3x+3\cdot 4)(x+1) & \text{\# раскрытие скобок} \\
                &= (x^2+1x+0)(x+1) & \text{\# вычисление в } \Z_6 \\
                &= x^3+x^2+x^2+x & \text{\# раскрытие скобок}\\
                &= x^3+2x^2+x
\end{align*}
Снова мы можем использовать Sage для вычислений с многочленами, коэффициенты которых находятся в $\Z_6$. \smecomb{(For the definition of rings see \ref{sec:rings}.)}{why is this ref here?} Для этого мы должны указать символ для неопределённой и тип для коэффициентов:
\begin{sagecommandline}
sage: Z6 = Integers(6)
sage: Z6x = Z6['x']
sage: Z6x
sage: p1 = Z6x([5,-4,2])
sage: p1
sage: p1 = Z6x([17,-4,2])
sage: p1
sage: Z6x(x-2)*Z6x(x+3)*Z6x(x-5) == Z6x(x^3 + 2*x^2 + x)
\end{sagecommandline}
\end{example}

Даны некоторые элементы того же типа, что и коэффициенты многочлена, многочлен может быть вычислен в этой точке, что означает, что мы подставляем данный элемент вместо каждого вхождения неопределённой $x$ в выражение многочлена.

Более точно, пусть $P\in R[x]$ с $P(x)=\sum_{j=0}^m a_j x^j$ — многочлен с коэффициентом типа $R$, и пусть $b\in R$ — элемент этого типа. Тогда \term{вычисление} (evaluation) $P$ в точке $b$ даётся следующим образом:
\begin{equation}
P(b) = \sum_{j=0}^m a_j b^j
\end{equation}
\begin{example}Рассмотрим снова многочлены с целыми коэффициентами из \examplename{} \ref{example:integer_polynomials}. Чтобы вычислить их в заданных точках, мы должны подставить точку вместо всех вхождений $x$ в выражение многочлена. Подставляя произвольные значения из $\Z$, мы получаем следующее:\sme{is this right?}
\begin{align*}
 &P_1(2)    = 2\cdot 2^2 -4\cdot 2 +17 = 17 \\
 &P_2(3)    = 3^{23}=94143178827 \\
 &P_3(-4)   = -4 \\
 &P_4(15)   = 174 \\
 &P_5(0)    = 1 \\
 &P_6(1274) =0 \\
 &P_7(-6)   = (-6-2)(-6+3)(-6-5) = -264 \\
\end{align*}
Заметим, однако, что невозможно вычислить любой из этих многочленов на значениях другого типа. Например, строго говоря, неверно писать $P_1(0.5)$, поскольку $0.5$ не является целым числом. Мы можем проверить наши вычисления с помощью Sage:
\begin{sagecommandline}
sage: Zx = ZZ['x']
sage: p1 = Zx([17,-4,2])
sage: p7 = Zx(x-2)*Zx(x+3)*Zx(x-5)
sage: p1(ZZ(2))
sage: p7(ZZ(-6)) == ZZ(-264)
\end{sagecommandline}

\end{example}
\begin{example} Рассмотрим снова многочлены с коэффициентами в $\Z_6$ из \examplename{} \ref{example:integer_mod_6_polynomials}\sme{check reference}. Чтобы вычислить их в заданных значениях из $\Z_6$, мы должны подставить точку вместо всех вхождений $x$ в выражение многочлена. Мы получаем следующее:
\begin{align*}
 & P_1(2)= 2\cdot 2^2 -4\cdot 2 +5 = 2 - 2 + 5 = 5\\
 &P_2(3)= 3^{23}=3\\
 &P_3(-4)= P_3(2) = 2\\
 &P_5(0)= 1\\
 &P_6(4)=0
\end{align*}
\begin{sagecommandline}
sage: Z6 = Integers(6)
sage: Z6x = Z6['x']
sage: p1 = Z6x([5,-4,2])
sage: p1(Z6(2)) == Z6(5)
\end{sagecommandline}

\end{example}
\begin{exercise}
Сравните оба раскрытия $P_7$ из $\Z[x]$ в \examplename{} \ref{example:integer_polynomials} и из $\Z_6[x]$ в \examplename{} \ref{example:integer_mod_6_polynomials}, и рассмотрите определение $\Z_6$, данное в \examplename{} \ref{def_residue_ring_z_6}. Можете ли вы увидеть, как определение $P_7$ над $\Z$ проецируется на определение над $\Z_6$, если рассматривать классы вычетов $\Z_6$?\sme{the task could be defined more clearly}
\end{exercise}
\subsection{Арифметика многочленов}\sme{this is the same as the section title}
Многочлены во многом ведут себя как целые числа. В частности, их можно складывать, вычитать и умножать. Кроме того, у них есть собственное понятие \concept{деления с остатком}. Неформально говоря, мы можем сложить два многочлена, просто сложив коэффициенты с одинаковыми индексами, а умножить их, применяя дистрибутивное свойство, то есть перемножив каждый член левого сомножителя с каждым членом правого сомножителя и сложив результаты.

Более точно, пусть $ \sum _{n = 0} ^{m_1}{a} _{n}{x} ^{n} $ и
$ \sum _{n = 0} ^{m_2}{b} _{n}{x^n} $ — два многочлена из $ R[x]$. Тогда \term{сумма} (sum) и \term{произведение} (product) этих многочленов определяются следующим образом:
\begin{equation}
\label{def:polynomial_arithmetic}
\sum _{n = 0} ^{m_1}{a} _{n}{x} ^{n} + \sum _{n = 0} ^{m_2}{b} _{n}{x } ^{n} = \sum _{n = 0} ^{max(\{m_1,m_2\})}{({a} _{n} +{b} _{n})}{x} ^{n}
\end{equation}
\begin{equation}
\label{def:polynomial_arithmetic_mul}
\bigg (\sum _{n = 0} ^{m_1}{a} _{n}{x} ^{n} \bigg) \cdot \bigg (\sum _{n = 0} ^{m_2 }{b} _{n}{x} ^{n} \bigg) = \sum _{n = 0} ^{m_1+m_2} \sum_{i = 0} ^{n}{a} _{i }{{b} _{n-i}}{x} ^{n}
\end{equation}
Правило вычитания многочленов может быть выведено из этих двух правил путём сначала умножения \uterm{вычитаемого} (subtrahend) на (многочлен) $-1$, а затем сложения результата с \uterm{уменьшаемым} (minuend).

Что касается определения степени многочлена, мы видим, что степень суммы всегда равна максимуму степеней обоих слагаемых, а степень произведения всегда равна степени суммы сомножителей, поскольку мы определили $-\infty + m= - \infty$ для каждого целого числа $m\in\Z$.
\begin{example}
\label{example:integer_polynomial_arithmetic_1}
 Чтобы привести пример того, как работает арифметика многочленов, рассмотрим следующие два многочлена с целыми коэффициентами $P,Q\in \Z[x]$ с $P(x)= 5x^2 -4x +2$ и $Q(x)=x^3-2x^2 +5$. Сумма этих двух многочленов вычисляется путём сложения коэффициентов каждого члена с равным показателем степени $x$. Это даёт следующее:
\begin{align*}
(P+Q)(x) & = (0+1)x^3 + (5-2)x^2 + (-4 +0) x +(2+5) \\
         & = x^3 +3x^2 -4x +7
\end{align*}
Произведение этих двух многочленов вычисляется путём перемножения каждого члена в первом сомножителе с каждым членом во втором сомножителе. Мы получаем следующее:
\begin{align*}
(P\cdot Q)(x) & = (5x^2 -4x +2)\cdot (x^3-2x^2 +5) \\
              & = (5 x^5 -10 x^4 +25 x^2)+ (-4x^4 +8 x^3 -20x) + (2x^3 -4x^2+10) \\
              & = 5 x^5 -14x^4 +10x^3+21x^2-20x +10
\end{align*}
\begin{sagecommandline}
sage: Zx = ZZ['x']
sage: P = Zx([2,-4,5])
sage: Q = Zx([5,0,-2,1])
sage: P+Q == Zx(x^3 +3*x^2 -4*x +7)
sage: P*Q == Zx(5*x^5 -14*x^4 +10*x^3+21*x^2-20*x +10)
\end{sagecommandline}
\end{example}
\begin{example}
\label{example:integer_mod_6_polynomial_arithmetic_1} Рассмотрим многочлены предыдущего \examplename{} \ref{example:integer_polynomial_arithmetic_1}, но интерпретированные в арифметике по модулю $6$. Так что мы снова рассматриваем $P,Q\in \Z_6[x]$ с $P(x)= 5x^2 -4x +2$ и $Q(x)=x^3-2x^2 +5$. На этот раз мы получаем следующее:
\begin{align*}
(P+Q)(x) & = (0+1)x^3 + (5-2)x^2 + (-4 +0) x +(2+5) \\
         & = (0+1)x^3 + (5+4)x^2 + (2 +0) x +(2+5) \\
         & = x^3 +3x^2 +2x +1\\
         \\
(P\cdot Q)(x) & = (5x^2 -4x +2)\cdot (x^3-2x^2 +5) \\
              & = (5x^2 +2x +2)\cdot (x^3+4x^2 +5) \\
              & = (5 x^5 +2 x^4 +1x^2)+ (2x^4 +2x^3 +4x) + (2x^3 +2x^2+4) \\
              & = 5 x^5 +4x^4 +4x^3+3x^2+4x +4
\end{align*}
\begin{sagecommandline}
sage: Z6x = Integers(6)['x']
sage: P = Z6x([2,-4,5])
sage: Q = Z6x([5,0,-2,1])
sage: P+Q == Z6x(x^3 +3*x^2 +2*x +1)
sage: P*Q == Z6x(5*x^5 +4*x^4 +4*x^3+3*x^2+4*x +4)
\end{sagecommandline}
\end{example}
\begin{exercise}
Сравните сумму $P+Q$ и произведение $P\cdot Q$ из двух предыдущих примеров \ref{example:integer_polynomial_arithmetic_1} и \ref{example:integer_mod_6_polynomial_arithmetic_1}, и рассмотрите определение $\Z_6$, данное в \examplename{} \ref{def_residue_ring_z_6}. Как мы можем вывести вычисления в $\Z_6[x]$ из вычислений в $\Z[x]$?
\end{exercise}
\subsection{\concept{Деление с остатком} многочленов}
Арифметика многочленов разделяет множество свойств с арифметикой целых чисел. Как следствие, понятие \concept{деления с остатком} и алгоритм деления столбиком также определены для многочленов. Вспомнив \concept{деление с остатком} целых чисел \ref{Euclidean_division}, мы знаем, что, данные два целых числа $a$ и $b\neq 0$, всегда существует другое целое число $m$ и натуральное число $r$ с $r<|b|$, такие что $a = m\cdot b +r$ выполняется.

Мы можем обобщить это на многочлены, когда старший коэффициент делимого многочлена имеет понятие мультипликативно обратного элемента. Действительно, данные два многочлена $A$ и $B\neq 0$ из $R[x]$, такие что $Lc(B)^{-1}$ существует в $R$, существуют два многочлена $Q$ (частное) и $P$ (остаток), такие что выполняется следующее уравнение и $deg(P) < deg(B)$:
\begin{equation}
\label{eq_polynomial_euclidean_division_notation}
	A = Q\cdot B + P
\end{equation}

Аналогично целочисленному \concept{делению с остатком}, и $Q$, и $P$ однозначно определены этими соотношениями.

\begin{notation}
\label{notation_polynomial_euclidean_division_notation}
Предположим, что многочлены $ A, B, Q $ и $ P $ удовлетворяют уравнению \ref{eq_polynomial_euclidean_division_notation}. Мы часто используем следующие обозначения для описания многочленов частного и остатка \concept{деления с остатком}:\footnote{\concept{Деление с остатком} многочленов объясняется более подробно в \cite{mignotte-1992}. Подробное описание связанного алгоритма можно найти в \chaptname{} 3, \secname{} 1 \cite{cohen-2010}. }
\begin{equation}
\label{def_polynomial_division_and_modulus}
\begin{array}{lcr}
\Zdiv{A}{B}: = Q, & & \Zmod{A}{B}: = P
\end{array}
\end{equation}
Мы также говорим, что многочлен $ A $ делится на другой многочлен $ B $, если выполняется $ \Zmod{A}{B} = 0 $. В этом случае мы также пишем $ B | A $ и называем $B$ \textit{множителем} (factor) $A$.
\end{notation}
Аналогично целым числам, методы вычисления \concept{деления с остатком} для многочленов называются \term{алгоритмами деления многочленов} (polynomial division algorithms). Вероятно, наиболее известным алгоритмом является так называемое \term{деление многочленов столбиком} (polynomial long division) (\algname{} \ref{alg_polynom_euclid_alg} ниже).
\begin{algorithm}\caption{Алгоритм Евклида для многочленов}
\label{alg_polynom_euclid_alg}
\begin{algorithmic}[0]
\Require $A,B \in R[x]$ with $B\neq 0$, such that $Lc(B)^{-1}$ exists in $R$
\Procedure{Poly-Long-Division}{$A,B$}
\State $Q \gets 0$
\State $P \gets A$
\State $d \gets deg(B)$
\State $c \gets Lc(B)$
\While{$ deg(P) \geq d$}
\State  $S := Lc(P)\cdot c^{-1}\cdot x^{deg(P)-d}$
\State $Q \gets Q + S$
\State $P \gets P - S\cdot B$
\EndWhile
\State \textbf{return} $(Q, P)$
\EndProcedure
\Ensure $ A=  Q \cdot B + P$
\end{algorithmic}
\end{algorithm}

% https://math.stackexchange.com/questions/2140378/division-algorithm-for-polynomials-in-rx-where-r-is-a-commutative-ring-with-u
Этот алгоритм работает только когда существует понятие деления на старший коэффициент $B$. Его можно обобщить, но нам понадобится лишь этот несколько более простой метод в дальнейшем.
\begin{example}[Деление многочленов столбиком] Чтобы привести пример того, как работает предыдущий алгоритм, разделим многочлен с целыми коэффициентами $A(x)=x^5+2x^3-9\in \Z[x]$ на многочлен с целыми коэффициентами $B(x)=x^2+4x-1\in\Z[x]$. Поскольку $B$ не является нулевым многочленом, а старший коэффициент $B$ равен $1$, который обратим как целое число, мы можем применить алгоритм \ref{alg_polynom_euclid_alg}. Наша цель состоит в том, чтобы найти решения уравнения \ref{eq_polynomial_euclidean_division_notation}\sme{check reference}, то есть нам нужно найти многочлен частного $Q\in\Z[x]$ и многочлен остатка $P \in \Z[x]$, такие что $x^5+2x^3-9 = Q(x)\cdot (x^2+4x-1) + P(x)$. Используя обозначение деления столбиком, которое в основном используется в англоязычных странах, мы вычисляем следующим образом:
\begin{equation}
\polylongdiv{X^5+2X^3-9}{X^2+4X-1}
\end{equation}
Таким образом, мы получаем $Q(x)=x^3-4x^2+19x-80$ и $P(x)=339x-89$, и, действительно, уравнение $ A=  Q \cdot B + P$ выполняется с этими значениями, поскольку $x^5+2x^3-9 = (x^3-4x^2+19x-80)\cdot (x^2+4x-1) + (339x-89)$. Мы можем проверить это, используя Sage:
\begin{sagecommandline}
sage: Zx = ZZ['x']
sage: A = Zx([-9,0,0,2,0,1])
sage: B = Zx([-1,4,1])
sage: Q = Zx([-80,19,-4,1])
sage: P = Zx([-89,339])
sage: A == Q*B + P
\end{sagecommandline}
\end{example}
\begin{example} В предыдущем примере деление многочленов дало нетривиальный (неисчезающий, то есть ненулевой) остаток. Деления, которые не дают остатка, представляют особый интерес. В этих случаях делители называются \term{множителями делимого} (factors of the dividend).

Например, рассмотрим снова многочлен с целыми коэффициентами $P_7$ из \examplename{} \ref{example:integer_polynomials}. Как мы показали, его можно записать как $x^3 - 4 x^2 - 11 x + 30$, а также как $(x-2)(x + 3)(x-5)$. Из этого мы видим, что многочлены $F_1(x)=(x-2)$, $F_2(x)=(x+3)$ и $F_3(x)=(x-5)$ являются всеми множителями $x^3 - 4 x^2 - 11 x + 30$, поскольку деление $P_7$ на любой из этих множителей даст нулевой остаток.
\end{example}
\begin{exercise} Рассмотрим выражения многочленов $A(x):= -3x^4 + 4x^3 + 2x^2 +4$ и $B(x)= x^2-4x+2$. Вычислите \concept{деление с остатком} $A$ на $B$ в следующих типах:
\begin{enumerate}
\item $A,B\in \Z[x]$
\item $A,B\in \Z_6[x]$
\item $A,B\in \Z_5[x]$
\end{enumerate}
Теперь рассмотрите результат в $\Z[x]$ и в $\Z_6[x]$. Как мы можем вычислить результат в $\Z_6[x]$ из результата в $\Z[x]$?
\end{exercise}
\begin{exercise}
Покажите, что многочлен $B(x)= 2x^4-3x+4\in \Z_5[x]$ является множителем многочлена $A(x)=x^7 + 4x^6 + 4x^5 + x^3 + 2x^2 + 2x + 3\in\Z_5[x]$, то есть покажите, что $B|A$. Чему равно $\Zdiv{A}{B}$?
\end{exercise}
\subsection{Простые множители} Вспомним, что основная теорема арифметики \ref{def:fundamental_theorem_arithmetic} говорит нам, что каждое натуральное число является произведением простых чисел. В этой главе мы увидим, что нечто подобное выполняется и для одночленных многочленов $R[x]$.\footnote{Строго говоря, это неверно для многочленов над произвольными типами $R$. Однако в этой книге мы предполагаем, что $R$ является так называемым факториальным кольцом, для которого содержание этого раздела справедливо.}
\footnote{Более подробное описание можно найти в \chaptname{} 3, \secname{} 4 \cite{mignotte-1992}.}


Аналогом простого числа для многочлена является так называемый \term{неприводимый многочлен} (irreducible polynomial), который определяется как многочлен, который нельзя разложить на произведение двух непостоянных многочленов с помощью \concept{деления с остатком}. Неприводимые многочлены для многочленов — это то же самое, что простые числа для целых чисел: они являются базовыми строительными блоками, из которых можно построить все остальные многочлены. 

Более точно, пусть $P \in R[x]$ — любой многочлен. Тогда всегда существуют неприводимые многочлены $F_1, F_2, \ldots, F_k \in R[x]$, такие что выполняется следующее:
\begin{equation}
\label{def_polynomial_prime_factorization}
P = F_1 \cdot F_2 \cdot \ldots \cdot F_k \;.
\end{equation}
Это представление единственно (за исключением перестановок множителей и умножения на ненулевые скалярные множители). Более того, каждый множитель $F_i$ называется \term{простым множителем} (prime factor) $P$.

\begin{example}
\label{example:irreducible_integer_polynomial_1}
 Рассмотрим выражение многочлена $P=x^2-3$. Когда мы интерпретируем $P$ как многочлен с целыми коэффициентами $P\in\Z[x]$, мы находим, что этот многочлен неприводим, поскольку любое разложение, отличное от $1\cdot(x^2-3)$, должно выглядеть как $(x-a)(x+a)$ для некоторого целого числа $a$, но не существует целых чисел $a$ с $a^2=3$.
\begin{sagecommandline}
sage: Zx = ZZ['x']
sage: p = Zx(x^2-3)
sage: p.factor()
\end{sagecommandline}
С другой стороны, интерпретируя $P$ как многочлен $P\in \Z_6[x]$ в арифметике по модулю $6$, мы видим, что $P$ имеет два множителя $F_1=(x-3)$ и $F_2=(x+3)$, поскольку
$(x-3)(x+3)= x^2 -3x +3x -3\cdot 3= x^2-3$.
\end{example}
Точки, в которых многочлен обращается в нуль, называются \term{корнями} (roots) многочлена. Более точно, пусть $P\in R[x]$ — многочлен. Тогда корень — это точка $x_0\in R$ с $P(x_0)=0$, а множество всех корней $P$ определяется следующим образом:
\begin{equation}
R_0(P):=\{x_0\in R\;|\; P(x_0)=0\}
\end{equation}
Корни многочлена представляют особый интерес относительно его разложения на простые множители, поскольку можно показать, что для любого данного корня $x_0$ многочлена $P$ многочлен $F(x)=(x-x_0)$ является простым множителем $P$. Корень $x_0$ многочлена $P$ называется имеющим \term{кратность} (multiplicity) $k$, если многочлен $(x-x_0)^k$ является множителем $P$, то есть если существует многочлен $Q$, такой что мы можем записать $P$ как
\begin{equation}
P(x)= (x-x_0)^k \cdot Q(x)
\end{equation}
Нахождение корней многочлена иногда называется \term{решением многочлена} (solving the polynomial). Это сложная проблема, которая была предметом многих исследований на протяжении всей истории.

Можно показать, что если $m$ — степень многочлена $P$, то $P$ не может иметь более $m$ корней. Однако в общем случае многочлены могут иметь меньше $m$ корней.
\begin{example}
Рассмотрим снова многочлен с целыми коэффициентами $P_7(x)=x^3 - 4 x^2 - 11 x + 30$ из \examplename{} \ref{example:integer_polynomials}. Мы знаем, что его множество корней даётся $R_0(P_7)=\{-3,2,5\}$.

С другой стороны, мы знаем из \examplename{} \ref{example:irreducible_integer_polynomial_1}, что многочлен с целыми коэффициентами $x^2-3$ неприводим. Отсюда следует, что он не имеет корней, поскольку каждый корень определяет простой множитель.
\end{example}
\begin{example}Чтобы привести ещё один пример, рассмотрим многочлен с целыми коэффициентами
$P=x^7 + 3 x^6 + 3 x^5 + x^4 - x^3 - 3 x^2 - 3 x - 1$. Мы можем использовать Sage для вычисления корней и простых множителей $P$:
\begin{sagecommandline}
sage: Zx = ZZ['x']
sage: p = Zx(x^7 + 3*x^6 + 3*x^5 + x^4 - x^3 - 3*x^2 - 3*x - 1)
sage: p.roots()
sage: p.factor()
\end{sagecommandline}
Мы видим, что $P$ имеет корень $1$, и что соответствующий простой множитель $(x-1)$ встречается в $P$ один раз. Мы также видим, что $P$ имеет корень $-1$, где соответствующий простой множитель $(x+1)$ встречается в $P$ $4$ раза. Это даёт следующее разложение на простые множители:
$$
P= (x - 1)(x + 1)^4(x^2 + 1)
$$
\end{example}
\begin{exercise}
Покажите, что если сумма кратностей всех корней многочлена $P\in R[x]$ степени $deg(P)=m$ меньше $m$, то многочлен должен иметь простой множитель $F$ степени $deg(F)>1$.
\end{exercise}
\begin{exercise}
Рассмотрим многочлен $P=x^7 + 3 x^6 + 3 x^5 + x^4 - x^3 - 3 x^2 - 3 x - 1\in \Z_6[x]$. Вычислите множество всех корней $R_0(P)$, а затем вычислите разложение $P$ на простые множители.
\end{exercise}
\subsection{\concept{Интерполяция Лагранжа}}
Одним особенно полезным свойством многочленов является то, что многочлен степени $m$ полностью определяется на $m+1$ точках вычисления, что подразумевает, что мы можем однозначно вывести многочлен степени $m$ из множества $S$:
\begin{equation}
\label{def_lagrange_interpolation_set}
S= \{(x_0,y_0), (x_1,y_1),\ldots,(x_m,y_m)\;|\; x_i\neq x_j\text{ для всех индексов $ i $ и $ j $}\}
\end{equation}
Многочлены следовательно обладают свойством, что $m+1$ пар точек $(x_i,y_i)$ для $x_i\neq x_j$ достаточно, чтобы определить множество пар $(x,P(x))$ для всех $x\in R$. Это свойство «чуть больше, чем нужно» многочленов широко используется, включая в SNARKs. Поэтому нам нужно понять метод фактического вычисления многочлена из множества точек.

Если коэффициенты многочлена, который мы хотим найти, имеют понятие мультипликативно обратного элемента, всегда возможно найти такой многочлен, используя метод, называемый \term{\concept{интерполяцией Лагранжа}} (Lagrange interpolation), который работает следующим образом. Дано множество вида \ref{def_lagrange_interpolation_set}, многочлен $P$ степени $m$ с $P(x_i)=y_i$ для всех пар $(x_i,y_i)$ из $S$ даётся \algname{} \ref{alg_lagrange_interplation} ниже.

\begin{algorithm}\caption{Интерполяция Лагранжа}
\label{alg_lagrange_interplation}
\begin{algorithmic}[0]
\Require $R$ must have multiplicative inverses
\Require $S= \{(x_0,y_0), (x_1,y_1),\ldots,(x_m,y_m)\;|\; x_i,y_i\in R, x_i\neq x_j\text{ для всех индексов $ i $ и $ j $}\}$
\Procedure{Lagrange-Interpolation}{$S$}
\For{$j \in (0\ldots m)$}
\State  $l_j(x) \gets \Pi_{i=0;i\neq j}^{m}\frac{x-x_i}{x_j-x_i} = \frac{(x-x_0)}{(x_j-x_0)} \cdots \frac{(x-x_{j-1})}{(x_j-x_{j - 1})} \frac{(x-x_{j+1})}{(x_j-x_{j+1})} \cdots \frac{(x-x_m)}{(x_j-x_m)}$
\EndFor
\State $P\gets \sum_{j=0}^m y_j\cdot l_j$
\State \textbf{return} $P$
\EndProcedure
\Ensure $P\in R[x]$ with $deg(P)=m$
\Ensure $P(x_j)=y_j$ for all pairs $(x_j,y_j)\in S$
\end{algorithmic}
\end{algorithm}

\begin{example} Рассмотрим множество $S=\{(0,4),(-2,1),(2,3)\}$. Наша задача — вычислить многочлен степени $2$ в $\mathbb{Q}[x]$ с коэффициентами из множества рациональных чисел $\Q$. Поскольку $\mathbb{Q}$ имеет мультипликативно обратные элементы, мы можем использовать метод \concept{интерполяции Лагранжа} из алгоритма \ref{alg_lagrange_interplation} для вычисления многочлена:
\begin{align*}
l_0(x) & = \frac{x-x_1}{x_0-x_1}\cdot\frac{x-x_2}{x_0-x_2}
         = \frac{x+2}{0+2}\cdot\frac{x-2}{0-2}
         =  -\frac{(x+2)(x-2)}{4}\\
       & = -\frac{1}{4}(x^2-4)\\
l_1(x) & = \frac{x-x_0}{x_1-x_0}\cdot\frac{x-x_2}{x_1-x_2}
          = \frac{x-0}{-2-0}\cdot \frac{x-2}{-2-2}
          = \frac{x(x-2)}{8}\\
       & = \frac{1}{8}(x^2-2x)\\
l_2(x) & = \frac{x-x_0}{x_2-x_0}\cdot\frac{x-x_1}{x_2-x_1}
          = \frac{x-0}{2-0}\cdot\frac{x+2}{2+2}
          = \frac{x(x+2)}{8}\\
       & = \frac{1}{8}(x^2+2x)\\
P(x)   & =  4\cdot (-\frac{1}{4}(x^2-4)) + 1\cdot \frac{1}{8}(x^2-2x) + 3\cdot \frac{1}{8}(x^2+2x)\\
       & = -x^2+4 + \frac{1}{8}x^2-\frac{1}{4} x + \frac{3}{8}x^2+\frac{3}{4} x \\
       & = -\frac{1}{2} x^2 +\frac{1}{2} x + 4
\end{align*}
И, действительно, вычисление $P$ на $x$-значениях $S$ даёт правильные точки, поскольку $P(0)=4$, $P(-2)=1$ и $P(2)=3$. Sage подтверждает этот результат:
\begin{sagecommandline}
sage: Qx = QQ['x']
sage: S=[(0,4),(-2,1),(2,3)]
sage: Qx.lagrange_polynomial(S)
\end{sagecommandline}
\end{example}
\begin{example} Чтобы привести ещё один пример, более актуальный для тем этой книги, рассмотрим то же множество, что и в предыдущем примере, $S=\{(0,4),(-2,1),(2,3)\}$. На этот раз задача состоит в том, чтобы вычислить многочлен $P\in\Z_5[x]$ из этих данных. Поскольку мы знаем
из \examplename{} \ref{primfield_z_5}, что в $\Z_5$ существуют мультипликативно обратные элементы, алгоритм \ref{alg_lagrange_interplation} применим, и мы можем вычислить единственный многочлен степени 2 в $\Z_5[x]$ из $S$. Мы можем использовать справочные таблицы из \eqref{Z5_tables} для вычислений в $\Z_5$ и получаем следующее:
\begin{align*}
l_0(x) & = \frac{x-x_1}{x_0-x_1}\cdot\frac{x-x_2}{x_0-x_2}
         = \frac{x+2}{0+2}\cdot\frac{x-2}{0-2}
         =  \frac{(x+2)(x-2)}{-4}
         =  \frac{(x+2)(x+3)}{1}\\
       & =  x^2+1\\
l_1(x) & =  \frac{x-x_0}{x_1-x_0}\cdot\frac{x-x_2}{x_1-x_2}
         = \frac{x-0}{-2-0}\cdot \frac{x-2}{-2-2}
         = \frac{x}{3}\cdot \frac{x+3}{1}
         = 2(x^2+3x)\\
       & =  2x^2+x\\
l_2(x) & = \frac{x-x_0}{x_2-x_0}\cdot\frac{x-x_1}{x_2-x_1}
         = \frac{x-0}{2-0}\cdot\frac{x+2}{2+2}
         = \frac{x(x+2)}{3}
         = 2(x^2+2x)\\
       & = 2x^2+4x\\
P(x)   & = 4\cdot (x^2+1) + 1\cdot (2x^2+x) + 3\cdot (2x^2+4x) \\
       & = 4x^2+4 + 2x^2 +x + x^2+2x\\
       & = 2x^2 +3x +4
\end{align*}
И, действительно, вычисление $P$ на $x$-значениях $S$ даёт правильные точки, поскольку $P(0)=4$, $P(-2)=1$ и $P(2)=3$. Мы можем проверить наши результаты с помощью Sage:
\begin{sagecommandline}
sage: F5 = GF(5)
sage: F5x = F5['x']
sage: S=[(0,4),(-2,1),(2,3)]
sage: F5x.lagrange_polynomial(S)
\end{sagecommandline}
\end{example}
\begin{exercise}
Рассмотрим арифметику по модулю $5$ из \examplename{} \ref{primfield_z_5} и множество
$S=\{(0,0),(1,1),(2,2),(3,2)\}$. Найдите многочлен $P\in\Z_5[x]$, такой что $P(x_i)=y_i$ для всех $(x_i,y_i)\in S$.
\end{exercise}
\begin{exercise}
Рассмотрим множество $S$ из предыдущего примера. Почему невозможно применить алгоритм \ref{alg_lagrange_interplation}, чтобы построить многочлен $P\in \Z_6[x]$, такой что $P(x_i)=y_i$ для всех $(x_i,y_i)\in S$?
\end{exercise}
